% ============================================================================
% Schizophrenia for a mathematician — English edition, expanded.
% Compiles with xelatex. Fonts live in fonts/, figures alongside.
% The evidence-tag macros keep their original names (\tU \tS \tP \tD) and only
% the printed letter changes, so the body text translates without touching tags.
% ============================================================================
\documentclass[11pt,a4paper,landscape,twocolumn]{article}
\usepackage{fontspec}
\usepackage{polyglossia}
\setdefaultlanguage{english}\setotherlanguage{russian}
\newcommand{\optserif}{Path=fonts/,Extension=.ttf,UprightFont=*-Regular,BoldFont=*-Bold,ItalicFont=*-Italic,BoldItalicFont=*-BoldItalic,Ligatures=TeX}
\newcommand{\optsans}{Path=fonts/,Extension=.otf,UprightFont=*-Regular,BoldFont=*-ExtraBold,ItalicFont=*-Italic,BoldItalicFont=*-ExtraBoldItalic,Ligatures=TeX}
\newcommand{\optmono}{Path=fonts/,Extension=.ttf,UprightFont=*-Regular,BoldFont=*-Bold,ItalicFont=*-Italic,BoldItalicFont=*-BoldItalic,Scale=0.9,Ligatures=TeX}
\setmainfont{LiberationSerif}[\optserif]
\setsansfont{Inter}[\optsans]
\setmonofont{LiberationMono}[\optmono]
\usepackage[a4paper,landscape,left=1.8cm,right=1.8cm,top=1.7cm,bottom=1.9cm,columnsep=1.0cm]{geometry}
\usepackage{graphicx,xcolor,amsmath,array,booktabs,enumitem,titlesec,fancyhdr,caption}
\usepackage[most]{tcolorbox}
\usepackage[hidelinks]{hyperref}
\definecolor{cB}{HTML}{2A78D6}\definecolor{cO}{HTML}{C9501F}\definecolor{cA}{HTML}{12875C}\definecolor{cY}{HTML}{A87200}
\definecolor{cM}{HTML}{5A6472}
\definecolor{pA}{HTML}{EAF7F1}\definecolor{pY}{HTML}{FCF4DD}\definecolor{pB}{HTML}{EDF3FB}
\raggedright\setlength{\parindent}{0pt}\setlength{\parskip}{5pt}\emergencystretch=2em
\titleformat{\section}{\sffamily\Large\bfseries\color{cB}}{\thesection.}{0.5em}{}
\titleformat{\subsection}{\sffamily\large\bfseries}{}{0em}{}
\titlespacing*{\section}{0pt}{16pt}{6pt}\titlespacing*{\subsection}{0pt}{8pt}{2pt}
\setlist[itemize]{leftmargin=1.3em,itemsep=2pt,topsep=2pt,parsep=0pt}
\setlist[enumerate]{leftmargin=1.6em,itemsep=3pt,topsep=2pt,parsep=0pt}
\pagestyle{fancy}\fancyhf{}\renewcommand{\headrulewidth}{0pt}\fancyfoot[C]{\sffamily\small\color{cM}\thepage}
\captionsetup{font={small,sf},labelfont=bf,labelsep=period,justification=raggedright,singlelinecheck=false,skip=4pt}
% Evidence tags: T textbook, R read in a paper, V verified by search, O our own.
\newcommand{\tU}{\textsuperscript{\sffamily\bfseries\color{cM}T}}
\newcommand{\tS}{\textsuperscript{\sffamily\bfseries\color{cB}R}}
\newcommand{\tP}{\textsuperscript{\sffamily\bfseries\color{cA}V}}
\newcommand{\tD}{\textsuperscript{\sffamily\bfseries\color{cO}O}}
\newtcolorbox{note}{colback=pA,colframe=pA,boxrule=0pt,arc=2pt,left=8pt,right=8pt,top=5pt,bottom=5pt,before skip=8pt,after skip=8pt}
\newtcolorbox{ineq}{colback=pY,colframe=pY,boxrule=0pt,arc=2pt,left=8pt,right=8pt,top=5pt,bottom=5pt,before skip=8pt,after skip=8pt}
\newcommand{\fig}[3][\textwidth]{\begin{figure*}[t]\centering\includegraphics[width=#1]{#2}\caption{#3}\end{figure*}}
\newcommand{\figp}[3][\textwidth]{\begin{figure*}[p]\centering\includegraphics[width=#1]{#2}\caption{#3}\end{figure*}}
\renewcommand{\dbltopfraction}{0.9}\renewcommand{\dblfloatpagefraction}{0.75}\renewcommand{\textfraction}{0.05}
\newcommand{\Hp}{H\textsubscript{2}O\textsubscript{2}}
\newcolumntype{R}[1]{>{\raggedright\arraybackslash}p{#1\linewidth}}
\newcommand{\facts}{\par\vspace{3pt}{\sffamily\bfseries\small\color{cM}Supporting facts}\par\vspace{-2pt}}
\raggedbottom
\begin{document}
\twocolumn[{%
{\sffamily\bfseries\fontsize{24}{28}\selectfont\color{cB} Schizophrenia for a mathematician\par}
\vspace{4pt}
{\sffamily\large The whole picture, from molecule to symptom\par}
\vspace{2pt}
{\sffamily\small\color{cM} Companion brochure to a project on a model of damage to cortical inhibitory cells. Version 6.0 (English, expanded), 3 November 2026.\par}
\vspace{14pt}}]

This brochure is written for a reader who is at home with dynamical systems but has not worked in biology and has not read the literature on the subject. Its aim is to give a connected picture: what processes run inside a cell and inside a network of cells, on what time scales, and where in that picture the nonlinearity appears that allows the illness to begin abruptly. Detail is deliberately sparse: every statement is given at the level needed to understand the next one.

The sections are built alike. First the principle in plain words. Then a short list of supporting facts with provenance tags. At the end, a box ``For the mathematician'': which mathematical object describes the thing.

\subsection{How to read the tags}
\begin{tabular}{@{}l@{\ \ }p{0.9\linewidth}@{}}
\tU & textbook knowledge: standard, written from memory, not separately checked;\\
\tS & taken from a paper read by the project's authors (figures from working notes);\\
\tP & verified by search while preparing the brochure, from an abstract, preprint or review;\\
\tD & an assumption or a calculation of our own, not present in the literature in this form.\\
\end{tabular}

All figures are our own diagrams, drawn for this brochure.

\subsection{Contents}
{\small
\begin{tabular}{@{}r@{\ \ }l@{\quad}r@{\ \ }l@{}}
1 & The whole picture on one page & 23 & Additional results\\
2 & The illness: what has to be explained & 24 & Slow defence\\
3 & Time scales & 25 & Priming\\
4 & The cell in five minutes & 26 & The precursor: what produces it\\
5 & Neuron and synapse & 27 & What an episode consists of\\
6 & The cortex: excitation and inhibition & 28 & How many cells are in the population\\
7 & The NMDA receptor & 29 & How the rhythm breaks\\
8 & Oxidants and defence & 30 & Asymmetric defence\\
9 & The switch & 31 & Three phases: the rescued band\\
10 & Loops and the developmental window & 32 & Mean field\\
11 & The day and sleep & 33 & Burstiness of the load\\
12 & Noise & 34 & The dopamine loop: a latch\\
13 & Dopamine and psychosis & 35 & The capacity loop\\
14 & A neighbouring hypothesis: BDNF & 36 & The sleep switch solved on paper\\
15 & How this is known & 37 & How we compute now\\
16 & Evidence: humans and animals & 38 & Accuracy and verifiability\\
17 & How a model comes out of this & 39 & Three more traps, no absolute null\\
18 & The sleep switch & 40 & What the model predicts\\
19 & From one cell to a population & 41 & What we hold beyond dispute\\
20 & The cortex in the loop & 42 & Open questions\\
21 & The conservation law for the gaps & 43 & Glossary\\
22 & What the model has already failed to do & 44 & Sources\\
\end{tabular}}

Sections 1--17 are biology and have changed little between versions. Sections 18--23 describe how the model is built now and what it has already managed to forbid; sections 24--33 cover the slow defence and priming, the attribution of the precursor, the composition of an episode, the number of cells as a parameter, the breakdown of the rhythm, the three phases of the population and what mean field will and will not give. Sections 34--40 are new in this edition: the two feedback loops and their verdicts, the closed-form solution of the sleep switch, the method as it now stands, and the list of what the model predicts that can be measured.


% ===================== BODY =====================
% =====================================================================
\section{The whole picture on one page}
\fig[0.8\textwidth]{fig_overview.png}{The working hypothesis as a single chain. Grey arrows are cause and effect; orange arrows are feedback.}
The hypothesis the brochure is built around has six steps (fig.~1).
\begin{enumerate}
\item The cortex contains a small group of inhibitory cells, the \emph{PV cells}. They hold the excitation of the rest in check and set the fast rhythm of the network.
\item PV cells work at their limit: many inputs, high firing rate, a large energy turnover. A by-product of that turnover is a flux of oxidants.
\item Enzymes protect the cell from oxidants. Part of that defence is produced only while the cell's NMDA receptor --- its activity sensor --- is working properly.
\item If the receptor works weakly, through heredity or through the environment, there is less defence and there are more oxidants. Oxidants in turn suppress the receptor: the loop closes on itself.
\item An overloaded PV cell does not die; it \emph{falls silent}, producing less of its working proteins. Inhibition weakens, the cortex is disinhibited, and the load on the surviving cells rises: a second loop.
\item An over-excited hippocampus, through a chain of intermediate structures, raises the responsivity of the dopamine system --- the bridge to the positive symptoms.
\end{enumerate}
\begin{note}
Every arrow in this chain has been shown separately in experiment: in cell culture, in brain slices or in animals. The chain as a whole has never been. A single hypothesis is set out here --- weak NMDA-receptor function together with impaired oxidative defence of PV cells. Its competitors are not reviewed.
\end{note}
\begin{ineq}
A circle that amplifies itself is positive feedback. If its gain is large enough, the system acquires two stable states and a threshold between them, and the transition between them is abrupt rather than gradual. That is what makes a sudden onset possible at all.
\end{ineq}

% =====================================================================
\section{The illness: what has to be explained}
Schizophrenia is a chronic psychiatric disorder. The latest revision of the international classification (ICD-11, in force since 2022, with countries migrating to it gradually) describes it not as a set of ready-made forms but as a profile along several independent symptom axes. Two things matter for a model: the profile differs between patients, and the course has a characteristic shape --- a long hidden preparation, a comparatively abrupt transition, and episodes thereafter.
\facts
\begin{itemize}
\item An ICD-11 diagnosis requires at least two of seven symptom groups, present most of the time for at least a month. At least one must come from the four core groups: delusions (persistent false beliefs); hallucinations, most often voices; thought disorder; and the experience of being influenced or controlled from outside. The remaining three are negative symptoms (less emotion, less drive, less speech), disorganised behaviour, and motor disturbance.\tP
\item The former forms (paranoid, hebephrenic, catatonic, simple) are abolished in ICD-11. Instead, six axes are rated separately for each patient: positive, negative, depressive, manic, motor and cognitive symptoms (memory, attention), together with the type of course: first episode, recurrent episodes, continuous.\tP
\item Catatonia --- a state of severe motor disturbance --- is now a separate diagnosis.\tP
\item Alongside schizophrenia, the same group of ``primary psychotic disorders'' contains separate diagnoses in which part of the picture is absent or short-lived: schizoaffective, schizotypal, delusional and acute transient psychotic disorder.\tP
\item \emph{Psychosis} is a state with pronounced positive symptoms; it arrives in episodes.\tU
\item Lifetime prevalence is 4.0 per 1000 (median across studies); lifetime morbid risk is 7.2 per 1000, roughly one person in a hundred and forty.\tP
\item Onset most often falls in late adolescence and early adulthood.\tU
\item Heritability is about 79\,\%.\tP\ This is the share of the variance of risk in the population attributable to genes, not the probability of passing the illness to a child. There is no single ``schizophrenia gene'': hundreds of variants are known, each contributing little.\tU
\item The first episode is preceded by a \emph{prodrome} --- months to years of non-specific change. The distribution of its length has a heavy tail: in one study the mean was 17.6 months and the median 4.\tS
\item In 43.7\,\% of patients the first psychosis developed acutely, over days.\tS
\item Drugs (antipsychotics) attenuate the positive symptoms and act little on the negative and cognitive ones.\tU
\end{itemize}
\begin{note}
The figures on prevalence, heritability and course were obtained under earlier diagnostic criteria (ICD-9, ICD-10 and the American manuals); the boundaries of the diagnosis shifted with the move to ICD-11. A long preparation followed by an abrupt transition suggests a system with a threshold, but a gradual accumulation of change with no threshold at all would produce a similar picture if symptoms appear only above some level of damage. The shape of the course is an argument, not a proof.
\end{note}
\begin{ineq}
A diagnosis is a point in a space of six axes, not the value of a single variable. A model of one mechanism may claim to explain individual axes, but not the illness as a whole. The hypothesis of this brochure bears chiefly on the cognitive axis (through PV cells and cortical rhythms) and on the positive one (through dopamine); about the depressive and manic axes it says nothing.
\end{ineq}

% =====================================================================
\section{Time scales}
\figp[0.9\textwidth]{fig_times.png}{Characteristic times of the phenomena discussed in the brochure. The scale is logarithmic.}
The processes relevant to the illness span more than ten orders of magnitude in time, from the millisecond impulse of a neuron to the multi-year maturation of the cortex (fig.~2). To the slow processes, the fast ones look instantaneously equilibrated; to the fast ones, the slow look like constant parameters.
\facts
\begin{itemize}
\item A neuronal impulse lasts about a millisecond.\tU\ A hallucination is usually seconds; only 7.4\,\% of episodes last 10 minutes or longer.\tS
\item The amount of a protein in a cell changes over hours and days.\tU
\item The brain's capacity to synthesise dopamine rises over months and years, from prodrome to psychosis.\tS
\item External rhythms impose their own scales: the day (light, sleep and wakefulness) and the year (day length, temperature).\tU
\end{itemize}
\begin{ineq}
A separation of time scales allows the model to be built in layers: fast variables are eliminated as algebraic relations, slow ones are frozen as parameters. One or two slow variables will not cover a spread this wide; the project's model describes the middle of the scale --- minutes to days, events inside a single cell.
\end{ineq}

% =====================================================================
\section{The cell in five minutes}
A cell is a droplet of aqueous solution inside a shell (the \emph{membrane}). Almost all the work inside it is done by \emph{proteins}. Some accelerate chemical reactions (these are called \emph{enzymes}), some serve as gates in the membrane, and some receive signals from outside (\emph{receptors}).

A protein is built from a gene in two steps: a working copy of the gene is made (mRNA), and the protein is assembled from it. The whole process is called gene \emph{expression}. A finished protein lives for a limited time and is then degraded. The amount of any protein is therefore a balance of synthesis and decay, and the cell governs itself by changing the rate of synthesis: it switches genes on and off in response to signals.

Energy is obtained in the \emph{mitochondria}, where nutrients are oxidised by oxygen. The energy is stored in ATP molecules and spent mainly on maintaining ion concentration differences across the membrane.
\facts
\begin{itemize}
\item Protein lifetimes range from minutes to weeks.\tU
\item The rate of an enzymatic reaction grows with the concentration of the substrate and saturates: there is a limited amount of enzyme, and each molecule of it is occupied for a finite time.\tU
\item The signal from a receptor reaches a gene along a chain of intermediate molecules; hours pass between the signal and an appreciable change in the amount of a protein.\tU
\item Most quantities in a cell are held in place by negative feedback.\tU
\item The brain is about 2\,\% of body mass and consumes about 20\,\% of the oxygen.\tU
\end{itemize}
\begin{ineq}
The amount of a protein $p$ obeys $p' = s(t) - p/\tau$: a first-order element whose time constant is the protein's lifetime. An enzymatic reaction is a saturating nonlinearity of the form $v = V s/(K + s)$. Control of genes supplies the slow variables, the chemistry of small molecules the fast ones. The saturation of the defence is exactly the source of the nonlinearity on which section~9 is built.
\end{ineq}

% =====================================================================
\section{Neuron and synapse}
\fig{fig_neuron.png}{Neuron and synapse. On the left the whole cell, on the right the point of contact between two cells, enlarged.}
A neuron is a cell adapted to transmitting signals. It maintains an electrical voltage across its membrane. When the summed input exceeds a threshold, a short impulse runs along the output process. Between cells the signal is passed chemically: the impulse releases molecules of a \emph{transmitter} from the terminal, and these open receptors on the next cell (fig.~3).
\facts
\begin{itemize}
\item A neuron consists of a body, branching \emph{dendrites} (input) and a single \emph{axon} (output). The impulse is all-or-none; information is carried by the rate and the timing of impulses.\tU
\item The main excitatory transmitter of the cortex is \emph{glutamate}, the main inhibitory one \emph{GABA}. Excitation brings the cell closer to threshold, inhibition takes it further away.\tU
\item Glutamate has two principal receptors. \emph{AMPA} is fast, acting in milliseconds. \emph{NMDA} is slow (tens to hundreds of milliseconds); section~7 is about it.\tU
\item Calcium entering the cell during its work serves as a signal to switch genes on. In this way synaptic activity governs which proteins the cell produces.\tU
\end{itemize}
\begin{ineq}
A neuron is a threshold element, a synapse a connection with a sign and a weight. What matters more for what follows is something else: through calcium, fast electrical activity (milliseconds) governs the cell's slow chemistry (hours). This is the bridge between time scales.
\end{ineq}

% =====================================================================
\section{The cortex: excitation and inhibition}
\fig{fig_circuit.png}{The simplest circuit of the cortex: excitatory (pyramidal) cells excite an inhibitory cell, which inhibits them in return. On the right, three molecules discussed below.}
Most cortical neurons excite one another. Such a network is unstable by itself: activity would run away to its limit. It is held in place by inhibitory cells, which receive excitation from their neighbours and inhibit them in return (fig.~4). The fastest of these contain the protein \emph{parvalbumin} and are called PV cells.
\facts
\begin{itemize}
\item About 80\,\% of cortical neurons are excitatory (\emph{pyramidal}) and about 20\,\% inhibitory.\tU\ PV cells make up roughly 40\,\% of the inhibitory cells of the mouse cortex.\tU
\item A single PV cell receives 16--34 thousand synapses; inhibition reaches a pyramidal cell with a delay of no more than 2\,ms; firing rates reach a hundred impulses per second.\tS
\item PV cells have a high energy turnover.\tS
\item Parvalbumin binds calcium and removes it quickly after an impulse.\tU\ \emph{GAD67} is the enzyme that makes GABA: less GAD67 means weaker inhibition.\tU
\item PV cells are surrounded by a shell of extracellular material (the \emph{perineuronal net}). It protects the cell from oxidants and matures during adolescence.\tS
\item PV cells set the \emph{gamma rhythm} --- oscillations of cortical activity at 30--80\,Hz. In patients with schizophrenia this rhythm is disturbed.\tU
\item The cortex operates in a regime where excitation amplifies itself and is held in check only by inhibition.\tS
\end{itemize}
\begin{ineq}
A PV cell is a fast negative feedback without which the network is unstable, and at the same time its clock generator. The price of speed is energy turnover, and therefore a flux of oxidants (section~8). This is the most heavily loaded element of the circuit.
\end{ineq}

% =====================================================================
\section{The NMDA receptor}
The NMDA receptor opens only when two conditions coincide: the transmitter has arrived \emph{and} the cell is already excited. Calcium enters the cell through it. It therefore serves not only as an input but as an activity sensor: from its operation the cell ``learns'' that the network needs it, and maintains production of its proteins, both working and protective.
\facts
\begin{itemize}
\item Ketamine blocks the NMDA receptor. In healthy volunteers it produces states resembling the symptoms of schizophrenia and impairs attention and memory.\tP
\item A paradox: an excitatory receptor is blocked, yet cortical activity rises. The accepted explanation is that receptors on inhibitory cells suffer more, inhibition weakens, and pyramidal cells are \emph{disinhibited}.\tU
\item In cell culture ketamine halves the amount of GAD67 in about 4 hours and of parvalbumin in about 8. The effect is reversible and the cells do not die.\tS
\item In mice a single dose of ketamine leaves no trace, whereas two doses on two consecutive days cause a persistent loss of parvalbumin and GAD67, with slow and incomplete recovery by day 10.\tS
\item Blocking the receptor in newborn rodents halves, within a day, the activity of the enzyme that makes the cell's principal antioxidant.\tS\ Synaptic activity switches on genes for protective enzymes in neurons, among them the one that repairs damaged defence (section~9).\tP
\item Oxidants reversibly suppress the NMDA receptor itself.\tS
\end{itemize}
\begin{ineq}
Three facts compose a closed loop: the receptor maintains the defence; without defence oxidants rise; oxidants suppress the receptor. The two-dose experiment shows a threshold not only in the strength of the insult but in its \emph{duration}: the system has a memory at least a day long.
\end{ineq}

% =====================================================================
\section{Oxidants and defence}\label{sec:redox}
\figp{fig_redox.png}{Where the peroxide comes from and where it goes. Below: the cycle of the protective enzyme peroxiredoxin, with the side branch that switches it off.}
An oxidant is a molecule that takes electrons from other molecules. Taking an electron from a protein or from DNA damages it. Oxidants are formed unavoidably wherever a cell works with oxygen. The principal one for what follows is hydrogen peroxide \Hp: it is comparatively long-lived and passes through membranes. In small amounts peroxide is an ordinary signal inside the cell; in large amounts it is a poison. \emph{Oxidative stress} is the state in which the influx of oxidants exceeds the capacity to clear them.

The defence has two layers. The first consists of enzymes that decompose peroxide; the fastest of these are the \emph{peroxiredoxins}. The second is a pool of small molecules that donate electrons and thereby return the enzymes to their working state; the principal such pool is \emph{glutathione}.
\facts
\begin{itemize}
\item Sources: the mitochondria (a by-product of energy production)\tU\ and the enzyme Nox2, which produces oxidants deliberately. In neurons Nox2 is switched on by excitation of NMDA receptors; its product leaves the cell and acts on neighbouring ones.\tS
\item The usual concentration of peroxide in a cell is 1--10\,nM.\tS
\item Glutathione is a short molecule of three amino acids; a cell contains five to six orders of magnitude more of it than of peroxide.\tU\ Half the glutathione of the mouse brain turns over in 60--80 hours.\tP\ Its level in the mouse brain varies 1.2-fold over the day.\tS
\item Neurons obtain the raw material for glutathione from \emph{astrocytes}, the brain's supporting cells.\tU
\item In mature neurons the system that switches on protective genes in response to stress itself is much weaker than in astrocytes.\tS
\end{itemize}
\begin{ineq}
The peroxide balance is fast: influx equals efflux within seconds, so the peroxide concentration is an algebraic function of the influx and the state of the defence. The nonlinearity arises because the efflux saturates and, worse, because the defence itself is damaged by excess peroxide (next section).
\end{ineq}

% =====================================================================
\section{The switch}
\fig{fig_switch.png}{Left: the steady fraction of the switched-off enzyme as a function of the peroxide influx (computed from the simplified equation; parameters chosen for clarity). Right: the same picture as a surface.}
Peroxiredoxin neutralises a molecule of peroxide by being oxidised itself, and returns to its working state within seconds. But if a second molecule of peroxide reaches the already oxidised enzyme in time, the enzyme is \emph{hyperoxidised} and switched off for a long time. It is repaired by a separate enzyme, \emph{sulfiredoxin}, and slowly.

Hence a positive feedback: the more peroxide, the more enzyme is switched off; the less working enzyme, the worse the clearance and the more peroxide. If the coupling is strong enough, the cell acquires two stable states at one and the same load: ``defence working'' and ``defence off'' (fig.~6).
\facts
\begin{itemize}
\item Repair of the hyperoxidised enzyme takes from 3 minutes to an hour in human cells.\tS
\item In mitochondria repair works differently: there is almost no sulfiredoxin there at rest. It is imported after peroxide leaves the mitochondria, reaches a maximum in 4--8 hours and is degraded once repair is complete.\tS
\item A detailed model of this system (Selvaggio et al., 2018; parameters measured for 15 cell types) yields two stable states for only one cell line and in a narrow range of loads.\tS
\item Reducing that model to a single equation gives an exact criterion for two states. The measured cells lie at the very boundary, but outside it: to move inside, repair must be weakened by a factor of 1.6--22.\tD
\item The threshold for cell death is sharp: doubling the peroxide dose takes cells from full survival to almost complete death.\tS
\end{itemize}
\begin{note}
Ordinary cells are \emph{not} in the region of two states. For the switch to work, the defence must be weakened. The hypothesis asserts that this is the case in PV cells with a weak NMDA receptor; there are no direct measurements for these cells. The kinetics of the enzyme have been measured in cell lines, not in neurons.
\end{note}
\begin{ineq}
This is a pair of folds (saddle-node bifurcations) with hysteresis: collapse and recovery occur at different loads. Inside the window a transition is possible only under a perturbation --- a push or noise (section~12). The width of the window and the height of the barrier are the quantities that have to be estimated quantitatively.
\end{ineq}

% =====================================================================
\section{Loops and the developmental window}
\fig[0.7\textwidth]{fig_loops.png}{The positive-feedback loops described in reviews, and the place of adolescence on the scale of a lifetime.}
The loop of section~7 is not the only one. Reviews describe several circuits passing through oxidative stress, and all of them converge on the PV cell (fig.~7). A separate question is why the disease begins in youth. The biologists' answer: PV cells and their protective shells mature late, towards the end of adolescence, and are vulnerable until then.
\facts
\begin{itemize}
\item Two principal loops between weak NMDA-receptor function and oxidative stress have been described: an intracellular one (receptor $\to$ protective genes) and a network one (disinhibition $\to$ inflammatory signal $\to$ Nox2 $\to$ oxidants).\tS
\item Another review counts five loops: through the receptor, the mitochondria, the brain's immune cells, the shells around the cells, and dopamine. The review contains no numbers.\tS
\item Genetic risk is usually compensated; the outcome depends on the strength \emph{and the duration} of the stress; mice with part of the receptor knocked out need a ``second hit'' to develop impairments.\tS
\item Oxidative stress damages PV cells specifically in childhood and adolescence.\tS
\item In rats a short course of a sedative in adolescence prevents subsequent impairments.\tS
\end{itemize}
\begin{ineq}
A ``vicious circle'' in a biological review is a set of arrows in a figure. Whether the loop closes depends on the product of the gains along it: if that product is less than one, a perturbation decays. Reviews do not pose this question, because they have no numbers. The developmental window is a slow drift of the defence parameters, that is, a slow passage through the boundary of the region of two states.
\end{ineq}

% =====================================================================
\section{The day and sleep}
\fig{fig_day.png}{What reaches its maximum and minimum, and when, over the day in rodents. Data from different studies, species and brain regions; for the bottom row the times are approximate.}
Almost every cell of the body runs a clock of its own with a period of about a day. Its mechanism is a negative-feedback loop with delay: two proteins switch on genes whose products, some hours later, suppress those same two proteins. The master clock sits in a small brain region above the crossing of the optic nerves; light through the eyes adjusts it to the alternation of day and night, and it synchronises the rest.

For our subject the day matters for two reasons. First, sleep and waking change neuronal activity and with it the influx of oxidants: the load on the defence is not constant, it passes through a cycle every day. Second, the oxidative defence itself oscillates with a daily period, and these oscillations have been found even in cells without a nucleus, that is, without genes.
\facts
\begin{itemize}
\item The clock's intrinsic period is preserved in constant darkness.\tU\ A daily rhythm of peroxiredoxin hyperoxidation is present in red blood cells, which have no nucleus.\tU
\item In the master clock of the mouse brain, the amount of hyperoxidised peroxiredoxin and the total amount of the enzyme both vary severalfold over the day; the hyperoxidised form leads the total amount by about 4 hours (figure reading; three animals).\tS
\item In slices of the master clock the tissue is reduced by day and oxidised at the beginning of the night; a change in the oxidative state alters neuronal excitability within about a minute.\tS
\item In PV cells of the rat frontal cortex a marker of oxidative DNA damage is maximal towards the end of the active phase and falls by 47\,\% by the end of the rest phase (fig.~8).\tP
\item Six hours of disrupted sleep raise this marker in rats by 25\,\%; in PV cells it is 27--50\,\% higher than in neighbouring cells.\tS\ After 4 hours without sleep the firing rate of cortical neurons is 34--64\,\% higher.\tS
\item In PV cells of the mouse frontal cortex 9\,\% of genes are read with a daily rhythm in males and 14\,\% in females (about 35\,\% in pyramidal cells); in males these are mainly genes for mitochondrial proteins. The intrinsic excitability of PV cells is the same by day and by night.\tP
\item In the mouse adrenal gland, heart and brown fat, mitochondrial hyperoxidised peroxiredoxin and the sulfiredoxin that repairs it oscillate with a daily period nearly in antiphase. Damage is maximal towards the end of the active phase, sulfiredoxin 8--12 hours later (figure reading; six points over the day). The total amount of peroxiredoxin itself is constant.\tS
\item Sleep and circadian disturbances occur in a substantial fraction of patients with schizophrenia.\tU
\end{itemize}
\begin{note}
Rodents are nocturnal: their active phase falls in darkness. Clock times cannot be carried over from mouse to human directly; what carries over is the phase relative to sleep and waking. The daily rhythm of hyperoxidised peroxiredoxin in the cortex and in PV cells has not been measured: the data are for the master clock and for organs outside the brain. In a whole-mouse-brain sample, mitochondrial hyperoxidised peroxiredoxin was not detected under ordinary conditions;\tS\ such a sample says nothing about a rare cell type.
\end{note}
\begin{ineq}
There are two distinct objects here. The first is a periodic external forcing of a system with two states: if the daily range of the load is comparable with the width of the window, the system approaches the boundary every day. The second is an intrinsic relaxation cycle ``damage $\to$ repair switched on $\to$ recovery $\to$ repair decays''; such a model has been published for cells outside the brain (del Olmo et al., 2019).\tS\ The two can be distinguished by the phase shift: a first-order element cannot lag its input by more than a quarter of the period, whereas the observed 8--12 hours out of 24 is between a third and a half.\tD
\end{ineq}

% =====================================================================
\section{Noise}
No quantity in the brain stands still. Neuronal activity fluctuates of its own accord, without external cause, and does so on all time scales at once: from fractions of a second to minutes and longer. The influx of oxidants fluctuates with it. For a system with two states this matters: a random push can carry it across the threshold before the mean load reaches the threshold.
\facts
\begin{itemize}
\item The spectrum of spontaneous brain activity has no preferred scale: the power decays with frequency as a power law over a range from about 0.003 to 500\,Hz.\tS
\item The level of arousal changes in jumps within seconds and returns to its previous value over tens of seconds.\tS
\item Episodes of ``voices'' in patients arise spontaneously, without external cause.\tS
\end{itemize}
\begin{ineq}
This is the problem of escape from a potential well under noise. The mean waiting time for a transition depends exponentially on the ratio of the barrier height to the noise intensity. Near the edge of the window the barrier decays as the distance to the edge to the power $3/2$.\tD\ The exponential dependence means that a modest spread of parameters between cells or between people produces an enormous spread of times --- a heavy-tailed distribution, as for the duration of the prodrome. A calculation of the project also revealed a limitation: fast noise of plausible strength is too weak to carry the system over far from the edge; what matters is slow fluctuation, on scales of minutes and longer.\tD
\end{ineq}

% =====================================================================
\section{Dopamine and psychosis}
\fig{fig_dopamine.png}{The path from the hippocampus to the dopamine cells, after the work of Grace's group. An arrow is excitation, a bar at the end inhibition.}
Damage to PV cells does not by itself explain why voices and delusions arise. The bridge to the symptoms is taken to be \emph{dopamine}, the transmitter with which the brain marks events as significant. If its release is excessive and disorderly, random events appear significant, and delusion arises as an attempt to explain them.
\facts
\begin{itemize}
\item All antipsychotics in use block dopamine receptors of one type (D2). The effect appears when about 65\,\% of the receptors are occupied; above 78\,\% motor side effects begin.\tP
\item In a synthesis of 44 studies, dopamine production and release are elevated in patients, while the number of receptors is almost unchanged.\tP
\item Substances that enhance dopamine release cause or exacerbate psychosis.\tU
\item Grace's scheme: weakening of PV cells in the hippocampus leaves it over-excited; through two inhibitory links this increases the fraction of dopamine cells ready to respond to a signal (fig.~9).\tS
\end{itemize}
\begin{ineq}
Damage to inhibitory cells changes the gain of the dopamine system. The project's model does not have this layer: it ends at the state of a single cell. Symptoms are a separate future problem, with a time scale of its own, of seconds.
\end{ineq}

% =====================================================================
\section{A neighbouring hypothesis: BDNF and serotonin}
It is not enough for a neuron to receive nourishment: to maintain its processes and synapses it needs support signals from its neighbours --- \emph{neurotrophic factors}. The principal one in the cortex is the protein BDNF. It is released by working excitatory cells and received by the TrkB receptor on neighbouring cells, PV cells among them. The \emph{neurotrophic hypothesis} of schizophrenia asserts that this support is insufficient: in development this yields faulty assembly of networks, and in the adult a reduced capacity of networks to reorganise.

\emph{Serotonin} is another transmitter which, alongside dopamine, governs the overall operating regime of the cortex. It has long been linked to psychosis: hallucinogens act on its receptors. One of these receptors governs BDNF production, so the two lines converge.
\facts
\begin{itemize}
\item In the frontal cortex of deceased patients the amount of BDNF mRNA is 28\,\% lower and that of its receptor TrkB 27\,\% lower. The more TrkB is reduced in a patient, the more GAD67 is reduced (correlation 0.74).\tP
\item In mice with the TrkB receptor reduced by 75\,\%, GAD67 and parvalbumin mRNA are 23 and 40\,\% lower. In mice with BDNF itself reduced by 80\,\% they are unchanged. The authors conclude that what matters is signalling through the receptor.\tP
\item In the blood of patients BDNF is moderately reduced (a synthesis of 41 studies, more than 7000 people). The reduction grows with the duration of the illness and is unrelated to symptom severity; whether it is cause or consequence the authors leave open.\tP
\item BDNF production depends on the cell's activity and on calcium, which enters in part through the NMDA receptor.\tU
\item The serotonin receptor 5-HT2A governs BDNF production in the frontal cortex and the hippocampus; the relation is bidirectional.\tP
\item Hallucinogens (LSD, psilocybin) act through the 5-HT2A receptor; most second-generation antipsychotics block not only the D2 dopamine receptors but this one as well.\tU
\item Antipsychotics and antidepressants raise BDNF production (per a 2005 review).\tP
\end{itemize}
\begin{note}
The neurotrophic hypothesis is not new: reviews under that name appeared in 2004--2005. It does not exclude the hypothesis of this brochure: both describe a weakening of PV cells, but place a different link first. Which of them is primary, and whether both are consequences of a common cause, has not been established. A search conducted while preparing the brochure found no single published model uniting serotonin, BDNF and oxidative stress.
\end{note}
\begin{ineq}
In terms of equations BDNF is another candidate for the role of the slow link ``network activity $\to$ support of the PV cell'': it is released by active cells and acts over hours and days. The construction is the same as in section~7: a slow variable which activity maintains and which itself maintains inhibition. The difference is that the source of support is the neighbouring excitatory cells, so the loop closes through the network rather than inside a single cell.\tD
\end{ineq}

% =====================================================================
\section{How this is known}
Almost everything said above was obtained not in sick people. The closer a method of measurement is to the molecule, the further it is from the disease; the reliability of a conclusion is determined by the rung on which it was obtained (table~1).
\begin{table*}[t]\small\centering
\begin{tabular}{@{}>{\raggedright\arraybackslash}p{0.17\textwidth}>{\raggedright\arraybackslash}p{0.37\textwidth}>{\raggedright\arraybackslash}p{0.38\textwidth}@{}}
\toprule
\textsf{\textbf{Rung}} & \textsf{\textbf{What it permits}} & \textsf{\textbf{What has to be paid}}\\
\midrule
Purified proteins and cell lines & measuring reaction rates; hence the kinetics of peroxiredoxin & these are not neurons; the conditions are far from a living cell\\
\addlinespace
Neuronal culture & following cells for hours, adding substances at an exact dose & the cells are taken from embryos or newborns, they are immature and deprived of their usual surroundings\\
\addlinespace
Brain slice & recording the current of an individual cell in an intact local network & a slice lives for hours; long-range connections are cut\\
\addlinespace
Living animal & intervening and observing the consequences over days and months; knocking out individual genes & a mouse does not have schizophrenia; a link is reproduced, not the disease\\
\addlinespace
Human: post-mortem tissue & counting cells and proteins in the region of interest & a terminal state after decades of illness and treatment\\
\addlinespace
Human: in-vivo imaging & measuring in patients and over time & an average over a block of tissue of cubic centimetres\\
\bottomrule
\end{tabular}
\caption{Rungs of evidence: from the molecule to the human being.}
\end{table*}
\facts
\begin{itemize}
\item \emph{Antibody staining}: a tissue slice is treated with a colour-labelled antibody against the protein of interest, and cells or brightness are counted. A cell that has ceased producing the labelled protein looks as though it has vanished. ``Loss of PV cells'' can therefore mean either death or silence.\tU
\item \emph{Separation of proteins in a gel}: proteins are separated by size, labelled with an antibody, and the brightness of the band is measured. The result is relative, usually as a fraction of a control or of a maximum.\tU
\item \emph{Recording a cell's current}: a fine electrode on the membrane of a single cell in a slice; it yields exact numbers for a few cells.\tU
\item \emph{Radioactive-label imaging} (PET) shows, for example, the capacity to synthesise dopamine; \emph{magnetic resonance spectroscopy} shows the concentration of a substance, glutathione in particular, averaged over all cells of the volume.\tU
\item Groups in animal experiments are small: three to ten animals per point is commonplace.\tU
\end{itemize}
\begin{ineq}
For model building three rules follow from this. A number measured on one rung cannot be substituted on another without qualification. Relative units yield ratios and phase shifts, but not absolute values. An average over a volume of tissue says nothing about a rare cell type: PV cells make up a few per cent of the neurons.
\end{ineq}

% =====================================================================
\section{Evidence: humans and animals}
\subsection{What is seen in patients}
\begin{itemize}
\item In the frontal cortex of deceased patients the amounts of parvalbumin and GAD67 mRNA are reduced, while the density of the PV cells themselves is unchanged: the cells are in place but produce less of these proteins.\tP
\item Glutathione. The first study (2000): 27\,\% less in the cerebrospinal fluid, 52\,\% less in the frontal cortex.\tS\ A 2021 synthesis (589 patients and 601 healthy subjects): the mean difference is not significant, but the spread in patients is larger.\tS
\item The capacity to synthesise dopamine is elevated (section~13).\tP
\end{itemize}
\subsection{What has been shown in experiments}
\begin{itemize}
\item \textbf{Ketamine in cell culture.}\tS\ GAD67 and parvalbumin fall by half in 4 and 8 hours; reversibly. This gives the time scale of the slow link and an argument for ``silence'' rather than death.
\item \textbf{Ketamine in mice.}\tS\ One dose, nothing; two consecutive days, a persistent loss. A threshold in duration.
\item \textbf{Receptor blockade in newborn rodents.}\tS\ Glutathione synthesis falls by half in a day. The link ``receptor $\to$ defence''.
\item \textbf{Prenatal brain injury in rats} (the substance MAM).\tS\ Loss of parvalbumin in the hippocampus, then excess dopamine activity in the adults. The bridge to dopamine and the role of adolescence.
\item \textbf{Mice lacking part of the NMDA receptor.}\tS\ A second hit is needed for impairments: a weak receptor shifts the system towards the boundary but does not take it across.
\end{itemize}
\begin{note}
The absence of a mean glutathione deficit in patients does not refute the hypothesis: spectroscopy averages over all cells, and a larger spread may mean a subgroup with a deep deficit or a disturbed daily rhythm. Nor does it confirm it: the explanation ``it is a matter of a subgroup or of a rare cell type'' requires independent verification, otherwise it renders the hypothesis irrefutable.
\end{note}

% =====================================================================
\section{How a model comes out of this}
\fig[0.86\textwidth]{fig_model.png}{The core of the model. A sharp arrow is enhancement, a bar at the end suppression. The colour indicates how well the link is founded.}
Ready-made computational models exist for the edges of the picture. The fast consequences for the network of a loss of parvalbumin and GAD67 have been modelled (Volman et al., 2011); so has the kinetics of peroxiredoxin (Selvaggio et al., 2018), as well as its daily oscillation outside the brain (del Olmo et al., 2019). No model of the damage loop itself --- receptor, oxidants, inhibitory cell --- was found in the literature by the authors of the project.\tS

The core of the project's model has two variables (fig.~10): a fast one $x$, the fraction of working peroxiredoxin, and a slow one $c$, the level of defence maintained by the work of the receptor. The peroxide concentration $u$ is obtained from an algebraic balance with the influx $\nu$:
\begin{align*}
&\nu = \varepsilon u + \frac{x\,u}{1+u},\quad \varepsilon = e_0 c^{\gamma},\qquad N = \frac{N_0}{1 + u/K_N},\\
&x' = c^{m}(1 - x) - \mu_0\,\frac{x\,u^2}{1+u},\\
&c' = \frac{1}{T}\Big(N + \frac{\beta u}{K_B + u} - c\Big).
\end{align*}
The first line and the equation for $x$ are derived from the measured kinetics. The equation for $c$ is an assumption about the functional form: the link ``receptor $\to$ defence'' itself is described in the literature. Time is measured in units of the enzyme's repair time; $T$ is the ratio of the lifetime of the protective proteins to it. The parameter $N_0 < 1$ represents a congenitally weak receptor.
\subsection{What has been computed}
\begin{itemize}
\item The slow link through the receptor widens the window of two states: the median ratio of its boundaries grows from 1.46 to 1.88; in 10.5\,\% of random parameter sets two states appear only because of it.\tD
\item A negative result. The core does not produce the asymmetry ``lost quickly, recovered slowly'' and does not reproduce the two-dose ketamine experiment: it has no memory between doses. Adding glutathione and feedback through network activity did not correct this.\tD
\item In 4.6\,\% of sets the core gives undamped ``damage--recovery'' oscillations with a period of about $3.4\,T$. Damage leads the defence by 0.16 of a period --- less than a quarter, as it must be for a first-order element; the shift observed for sulfiredoxin, from a third to a half of a period, is inaccessible to the core.\tD
\item With random fluctuations of the load, the law ``the barrier decays as the distance to the edge to the power $3/2$'' and the formula for the mean waiting time of a transition were confirmed by computation. The barriers turned out to be high for fast noise.\tD
\end{itemize}

% =====================================================================
\section{The sleep switch}\label{sec:sleep}
\fig{fig_sleepswitch.png}{Left: two nuclei inhibit each other, so two states are possible at one and the same input. The arrows are the transitions at the folds. Right: the ``sleep pressure'' accumulated over the day under one and the same schedule.}

The load on a cell is not constant: by day the brain works, by night it rests. For this to enter the model honestly, the sleep schedule must not be imposed by hand --- sleep has to come out by itself.

It works as follows. There is a ``sleep'' nucleus and there are wake nuclei; they inhibit each other. Such a pair behaves as a trigger: at one and the same external input it has two stable states, and the system stays in whichever one it finds itself. The transition happens not when some threshold is crossed, but when the branch the system was sitting on simply ceases to exist --- such a point is called a fold.

The input to the sleep nucleus is governed by two things: an internal clock with a period of about a day, and the ``sleep substances'' accumulated during waking. The latter grow while one is awake and decline during sleep. When they become sufficient, the upper branch disappears and one falls asleep; when over the night they become low, the lower branch disappears and one wakes.

\facts
\begin{itemize}
\item The model of the pair of nuclei and all nine of its numbers were taken from the paper in full and were not tuned. Check: our folds came out at 1.450 and 2.463\,mV, against 1.45 and 2.46 in the paper; the free-running cycle gave 15.5\,h of waking, 8.5\,h of sleep, a period of 24.0\,h.\tS
\item Time constants of the sleep substances: rise 18.2\,h, decay 4.2\,h. The decay is more than three times faster than the rise --- sleep restores faster than waking loads.\tS
\item In the original paper a single constant of 45\,h stood for both directions. The sawtooth is then almost straight and, more importantly, sleep restores at exactly the rate at which waking loads, which cannot be. After the substitution, the range of the sleep substances came out at 0.527 against 0.50 in the authors of the two-process model --- and this agreement was not fitted by anyone.\tD
\item The night-to-night spread of sleep duration in humans: 1.29\,h by actigraphy (26\,121 nights), 1.12\,h by EEG, 0.54\,h by self-report.\tP
\end{itemize}

\begin{ineq}
A mutually inhibiting pair yields an S-shaped curve of steady states; the middle branch is unstable, the outer ones stable. The transitions are saddle-node bifurcations at the ends of the middle branch. The slow variable (the sleep substances) creeps along the curve and drops the system off the fold; this is a relaxation oscillator entrained by an external periodic term. What matters is that the position of the folds is a property of the fast subsystem, not a threshold that could be assigned.
\end{ineq}

% =====================================================================
\section{From one cell to a population}\label{sec:pop}
\fig{fig_population.png}{A single cell has two states and nothing in between. Cells differ in resilience, so the fraction damaged changes smoothly --- but the hysteresis loop is preserved.}

So far the discussion has been about a single cell. But there are many inhibitory cells in the cortex, and a population behaves differently from any of its cells.

A single cell with two states is a switch: either healthy or damaged, with nothing intermediate. This is visible in the computation as well: over the whole run a single cell spent about one per cent of the time in the middle of the damage scale. That is, it has no slow deterioration at all: only a jump.

If the cells differ slightly in resilience --- and they do --- then under a growing load they break not all at once but in turn. The fraction broken grows smoothly. The hysteresis loop does not go anywhere: for the cells to return, the load must be lowered below the one at which they broke. The result is a smooth quantity that has retained memory --- something a single cell does not have.

There is a subtlety here which forbids simply averaging. If one takes a population of independent cells and computes the mean, smoothness comes out ``from nothing'' --- purely from averaging, without any mechanism. For this reason the cells in the model are coupled through what couples them in the cortex: through network activity. Fewer working inhibitory cells, weaker inhibition, higher excitation, higher load on those remaining. The circle closes in fact and not in words.

\facts
\begin{itemize}
\item The mean over a population of bistable units with scattered thresholds changes smoothly and preserves hysteresis.\tP
\item The intrinsic heterogeneity of the population does not collapse into a single mean-field equation: what is needed is the fraction broken, not the mean state.\tP
\item Computation: for a single cell 0.9--1.3\,\% of the time falls in the middle of the scale; for a population of 32 cells, about 5\,\%.\tD
\item A negative result. The smoothness is produced by the spread of resilience and not by the coupling through the network: with the narrowest spread it almost vanishes. The hope that the circle ``activity $\to$ load $\to$ damage $\to$ inhibition $\to$ activity'' would itself create intermediate states was not borne out.\tD
\end{itemize}

\begin{ineq}
Formally: a population of identical bistable elements has the same bifurcation diagram as a single element. A spread of a parameter smears both folds into an interval, and the fraction of elements in the upper state becomes the distribution function of the thresholds --- a step function with step $1/J$, monotone and hysteretic. A common coupling through the mean field adds a positive feedback to this, so it can narrow the interval back or even restore a discontinuous transition.
\end{ineq}

% =====================================================================
\section{The cortex in the loop}\label{sec:loop}
For the circle to close, a model of the cortex is needed --- a crude one at least. We take two groups of cells: excitatory, and inhibitory PV cells. Each is described by one number, the mean firing rate. The excitatory cells drive themselves and the inhibitory ones; the inhibitory cells damp both.

Damage enters here as a single quantity: the output weights of the inhibitory cells are multiplied by a factor which falls with the fraction damaged. At complete loss, inhibition drops to 60\,\% of its original value --- this number is taken from a study in which such a loss was modelled.

What this coupling gives: when inhibition falls to 60\,\%, the activity of the excitatory cells rises by 19\,\%. The load on the remaining inhibitory cells rises by the same amount. This is the closed circle, and its strength is not a fitted number but a consequence of two taken from the literature.

\facts
\begin{itemize}
\item Neural-mass models in which the inhibitory weight is itself a variable with an equation of its own exist in the literature; so does the fall of inhibition to 60\,\% of its original value under loss of PV cells.\tP
\item The factor is applied to all output weights of the inhibitory group at once rather than to one: this is how it was done in the study that analysed PV loss specifically.\tP
\item No models linking a cell's metabolism to a neural mass could be found. The circle ``over-excitation $\to$ damage to metabolism $\to$ less inhibition $\to$ over-excitation'' is described in the literature in words, without equations.\tS
\end{itemize}

\begin{note}
The cortex is taken here in a steady state: its own time scales are tens of milliseconds, that is, thousands of times faster than everything else. For this reason the gamma rhythm is deliberately absent from this version: what is needed is the mean firing rate, on which metabolism depends. The regime with oscillations is a separate problem.
\end{note}

% =====================================================================
\section{The conservation law for the gaps}\label{sec:law}
\fig{fig_gaps.png}{An episode begins when the load is thrown above the top of the window, and ends when it goes below the bottom. The sum of the two jumps does not depend on where in the window the mean load sits.}

This is the principal result the model produced not as a number but as a prohibition.

The picture of the disease that has to be obtained: episodes are rare and last from hours to several days. So a collapse must be hard and a return easy. In the model both are done by one and the same random fluctuation of the load: upwards it throws the system over the upper boundary of the window, downwards it returns it below the lower one.

And here a simple thing comes to light. The distance up and the distance down, measured in the logarithm of the load, sum to the width of the window --- and to nothing else. Shift the mean load down and the collapse becomes far away, but the return becomes near; shift it up and the reverse. The sum does not change.

For a typical parameter set this sum is about nine standard deviations of the noise. The picture of the disease requires three deviations upwards (so that collapses are rare) and one and a half downwards (so that the return takes a day) --- four and a half in all. Nine does not fit into four and a half.

The conclusion is hard and independent of the parameters: \textbf{as long as collapse and return are done by one and the same noise, there will be no rare episodes with a fast return.} Any numbers may be turned; it will not help.

What follows from this. An asymmetry is needed: collapse and return must be governed by different things. Two routes are already visible, and both are testable.

\begin{itemize}
\item \textbf{A one-sided trigger.} A serotonin surge raises the load and never lowers it. The upward gap is then covered by the surge, while the downward gap is left to the ordinary noise, and the two distances are measured by different distributions.
\item \textbf{Variable defence.} If after a load the cell strengthens its repair, the lower boundary of the window itself moves upwards, and the return occurs without a large excursion. The law then ceases to bind from the other end --- it was derived on the assumption that the boundaries of the window stand still.
\end{itemize}

\facts
\begin{itemize}
\item Activation of the 5-HT2A receptor raises the firing rate of cortical excitatory cells roughly twofold; the pathway is confirmed by the effect being abolished both by blockade of the receptor and by blockade of the intracellular link. This is a pharmacological maximum, not an ordinary surge.\tP
\item Duration of the acute action of substances working through this receptor: 6--7\,h and 10--11\,h, with onset after 0.4--0.8\,h.\tP
\item The peroxiredoxin repair enzyme is up-regulated by activity: mRNA 13.5-fold in 4 hours, the protein by one day. The rate of the reverse decay has not been measured in any of the papers read.\tS
\end{itemize}

\subsection{A refinement that was not borne out}
In version~3.1 an explanation stood here: that an episode fails to end because a random dip in the load is too short --- the cell repairs itself over about a day, whereas the fast component of the noise persists for two hours. Two predictions followed from it: the dependence of episode duration on the correlation time of the noise should have a minimum near the repair time, and lengthening the fast component should shorten the episode.

Verified in a separate run. \textbf{Both predictions proved false.} There is no minimum at all: the duration grows monotonically, and the shorter the noise the shorter the episode, all the way to the edge of the scan. Lengthening the fast component does not shorten the episode but lengthens it --- from 3.4 to 7.9 days. The explanation is withdrawn.

What remains in its place is a measured dependence, without an explanation:
\[\text{duration} \;\propto\; \tau^{1/4}\,\sigma^{2/3},\]
where $\tau$ is the correlation time of the noise and $\sigma$ its strength. The exponent of $\tau$ came out the same for both components, slow and fast: 0.254 and 0.243, over ranges of values of 40-fold and 33-fold. Such an agreement itself calls for an explanation, but to give one before verification would be a repetition of the same error.

A practical conclusion follows at once. At a correlation time of six hours the episodes come out at 1.5 days, and more than half of them are shorter than a day --- precisely the picture we were after. But the sleep precursor described in section~\ref{sec:supp} requires noise with a time of several days: otherwise there is simply nothing for it to arise from. Both things cannot be obtained from one source of randomness. Hence the episode must end not at the will of the noise but of itself --- and this is one more argument in favour of variable defence.

\begin{ineq}
Let $W$ be the ratio of the window's boundaries, $q$ the position of the mean load within it on a logarithmic scale, $F$ the feedback factor at complete damage, and $P$ and $T$ the daily peak and trough. Then
\[\underbrace{\ln\frac{W^{1-q}}{P}}_{\text{to collapse}} + \underbrace{\ln\big(W^{q} F T\big)}_{\text{to return}} = \ln\frac{W F T}{P},\]
and the right-hand side does not depend on $q$. The statement is elementary --- it is the addition of two logarithms --- but it constrains any model of this construction: a bistable unit under a noisy load cannot simultaneously have a high entry barrier and a low exit barrier. Hence the requirement of an asymmetric mechanism.
\end{ineq}

% =====================================================================
\section{What the model has already failed to do}\label{sec:fail}
The negative results are listed here separately and deliberately. Each of them is a check written down before the computation which did not pass; each narrowed the range of possible explanations more than any of the successes did.

\begin{enumerate}
\item \textbf{The daily range is not a gate but a volume knob.} It was expected that there exists a threshold value of the daily range of the load below which collapses are almost absent and above which they occur. What came out was a smooth growth with no break: over a range from 0 to 40\,\% the fraction of cells that were damaged grew only 1.5-fold. The range changes the risk but switches nothing on.\tD
\item \textbf{Smoothness comes from the spread, not from the network.} The hope was that coupling through network activity would itself create intermediate states between ``healthy'' and ``damaged''. It does not: with a narrow spread of cell resilience the smoothness almost vanishes (fig.~\ref{fig:res}a).\tD
\item \textbf{There is no sleep precursor without a slow variable.} As long as the only source of randomness varied over 2.4 hours, no change of sleep in the day before a collapse arose --- nor could it: the model contained no quantity living on the scale of days. The precursor appeared only after such a quantity was added.\tD
\item \textbf{The episode duration remained wrong.} In the model an episode lasts about six days, whereas observation gives from hours to a few. This is a direct consequence of the law of section~\ref{sec:law}, and it is not corrected by parameters.\tD
\item \textbf{Replacing the ``fraction broken'' by the ``mean damage'' changed nothing.} The argument was strong: the fraction is zero until the first cell breaks, whereas the mean grows at once. In fact, in a narrow window the cells break almost simultaneously, and the two quantities agreed to within half a per cent.\tD
\item \textbf{The serotonin surge turned out to be not a trigger but the same volume knob.} The check was direct: rare one-sided surges raising the load were to cover the upward gap and leave the return to the ordinary noise. Of the rule's four parts not one passed. Collapses following a surge amounted to 20\,\% instead of the expected 80; the episode duration did not change at all --- 5.8 days without surges and 5.2 with them. Surges add collapses (the fraction of affected sets grows from 0.59 to 0.73 as surges become more frequent) but change no property of the episode itself.\tD
\item \textbf{The episode does not remember how deep it went.} What was tested was whether an episode lives a life of its own after the load has returned to normal. It turned out that it almost does not: the episode duration coincides with the time the load spent above the upper boundary of the window, with a rank association of 0.94, whereas the association of the recovery tail with the depth of damage attained is only 0.11. The tail meanwhile occupies a quarter of the episode and lasts about half a day regardless of how many cells broke.\tD
\item \textbf{The first rule turned out to be almost empty.} The check on damage leading the defence passed, but it then emerged that a first-order element satisfies it almost always. A lead of a third of a period, which is observed in experiments, is inaccessible to this core for any parameters.\tD
\end{enumerate}

\begin{note}
Twice during the work a rule passed ``idly'': it was formally satisfied, but on examination it turned out that an extraneous effect was responsible. Once different sets were being compared instead of one and the same; the other time the control rested on a single set, and the statistic simply did not compute. Both cases came to light only because for every rule not only the check but also a way of fooling it had been written down in advance.
\end{note}

% =====================================================================
\section{Additional results}\label{sec:supp}
\figp{fig_results.png}{\label{fig:res}Four results that are visible without preparation. The numbers are medians over several hundred random parameter sets.}

The working arrangement is as follows: before anything is computed, a rule is formulated --- what exactly must come out, and with what number --- and only then is the computation launched. A rule may not be rewritten after the result has been seen; it may only be declared unfit, with a statement of why. Below are all the rules of the project and what came of them.

\begin{table*}[p]\small\centering
\begin{tabular}{@{}R{0.04}R{0.30}R{0.32}R{0.27}@{}}
\toprule
 & \textbf{what was checked} & \textbf{what came out} & \textbf{verdict}\\
\midrule
F1 & damage leads the defence by 1/8--1/4 of a period & 0.162 & satisfied, but the rule judged almost empty\\
F2a & collapses accumulate towards the second half of waking & 0.88 against 0.50 at random; analytic expectation 0.82 & satisfied, and in agreement with the derivation\\
F2b & there is a threshold in the daily range & 1.5-fold growth, no break & \textbf{not satisfied}\\
F2b$'$ & the curve lies on the smooth prediction without a gate & worst deviation 0.55 against a tolerance of 0.04 & \textbf{not satisfied}: neither a gate nor a smooth hazard\\
F2c & sleep rarely interrupts an episode & 0.51 against a uniform expectation of 0.35 & \textbf{not satisfied}: sleep interrupts \emph{more} often than at random\\
F3a & damage impedes falling asleep and lengthens the episode & 1.67-fold against a threshold of 1.5 & satisfied\\
F3a$'$ & \dots and by exactly as much as waking lengthens & observation 1.67 against 1.10 for the switch & unverifiable; the amplification is real and requires explanation\\
F3c & episodes shorter than a day are rare & 0.13 over all episodes against a threshold of 0.50 & satisfied\\
F4a & sleep deteriorates before the collapse threshold & 1.085 against a threshold of 1.10 and a switch expectation of 1.114 & \textbf{not satisfied}, and below the mechanical expectation\\
F4b$'$ & the precursor exists relative to each set's own level & 1.055 against a threshold of 1.05 & satisfied\\
F5 & the fraction damaged has become smooth & 0.038 against a threshold of 0.05 and a derivation bracket of 0.069--0.139 & \textbf{not satisfied, and below the bracket}: hysteresis removes the intermediate states\\
F6a & waking lengthens before the collapse & 1.055, inside the mechanical bracket & satisfied\\
F6b,d & the slope of the precursor curve & $-0.62$ over eight cells, MDE 0.94 & unverifiable at any sample size: the cells must be subdivided\\
F6c$'$ & slow noise raises the precursor (paired) & $p=0.97$ & \textbf{not satisfied}: the noise has nothing to do with it\\
F7a & the episode is no longer than two days & 4.39 days & \textbf{not satisfied}\\
F7a$'$ & the duration follows the law of F16 outside the fit & miss of 0.22 days against a tolerance of 1.5 & satisfied\\
F7b & there exist episodes shorter than a day & 0.28 over all episodes against a threshold of 0.30 & \textbf{not satisfied}\\
F8a & the collapse is caused by a serotonin surge & 0.24 against a threshold of 0.80 and a derivation expectation of 0.36 & \textbf{not satisfied}: the threshold was a requirement on the selection\\
F8b,c & the surge changes the properties of the episode & 7.2 days; short ones 0.12 & \textbf{not satisfied}\\
\bottomrule
\end{tabular}
\caption{Pre-registered rules of the project, part one: the core, the day, the episode and the serotonin surge.}
\end{table*}

\begin{table*}[p]\small\centering
\begin{tabular}{@{}R{0.04}R{0.30}R{0.32}R{0.27}@{}}
\toprule
 & \textbf{what was checked} & \textbf{what came out} & \textbf{verdict}\\
\midrule
F10a & the recovery tail is related to the depth of damage & rewritten for the present model, awaiting a run & written before the computation\\
F10b,c & the episode is the time of high load, not the repair & provisionally 0.87 and a tail of 7\,\% & written before the computation\\
F11 & induction closes the episode & the association of induction strength with duration is weaker than the threshold & \textbf{not satisfied}\\
F12 & the defence creates a refractory period & by the intervals between episodes, no & \textbf{not satisfied}\\
F13 & the response does not depend on the number of cells & ranges 0.21 and 0.24 against noise ceilings of 0.055 and 0.042 & \textbf{not satisfied}: $J$ is a parameter, not a resolution\\
F14 & the memory of the defence is masked by noise & two parts out of three & partially\\
F15 & a refractory period exists and is not due to the defence & $t_r=1.0\dots1.4$ days, both controls passed & satisfied\\
F16 & the duration law predicts outside the fit & miss of 10\,\% against a tolerance of 2\,SE & satisfied\\
F16c & the law is affine, not saturating & the straight line beat both saturating forms & satisfied\\
F18b$'$ & the precursor is made by the memory of the defence (paired) & falls in 107 sets, $p\le0.05$ & satisfied\\
F18b$''$ & \dots and makes it entirely & only a fifth of the excess goes away & \textbf{not satisfied}: there is a third slow quantity\\
F19 & priming brings the episode to the target & 1.46 days, bound 1.665, $p=5.2\cdot10^{-6}$ & \textbf{satisfied, all four parts}\\
F20 & the rhythm breaks at strong coupling & locking staircase 1:1 $\to$ 5:4 $\to$ 3:2 $\to$ 2:1 $\to$ 3:1 $\to$ 4:1 $\to$ no sleep & satisfied; this is a staircase, not a period-doubling cascade\\
F20e & there exist regimes with alternating cycle lengths & patterns 25,\,24,\,24,\,47 and 47,\,25 and 46,\,26,\,24\,h & satisfied; the working region does not reach them\\
F21 & the third slow quantity is the filling of the population & written before the computation & awaiting a run\\
F22 & $J$ enters through the effective spread of the grid & the arithmetic of the grid confirmed; the check by substituting sigma written down & awaiting a run\\\\
\bottomrule
\end{tabular}
\caption{Pre-registered rules of the project, part two: the tail, the refractory period, the duration law, the slow defence, priming and the population.}
\end{table*}

\subsection{Comparison with published numbers}
Separately from the rules there are checks in which the model is obliged to reproduce someone else's result. This is not a success of the model but a condition of admission to work: if it does not agree, no further computation is permitted.

{\small
\begin{tabular}{@{}R{0.34}R{0.28}R{0.30}@{}}
\toprule
\textbf{quantity} & \textbf{our computation} & \textbf{in the paper}\\
\midrule
folds of the sleep switch & 1.450 and 2.463\,mV & 1.45 and 2.46\,mV\\
waking / sleep / period & 15.50 / 8.50 / 24.00\,h & 15.5 / 8.5 / 24.0\,h\\
range of the sleep pressure & 0.527 & 0.50\\
spread of sleep duration & 0.67\,h & 1.12--1.29\,h (with the daily routine)\\
integration step against an exact solver & discrepancy $3\cdot10^{-7}$ & ---\\
\bottomrule
\end{tabular}}

The range of the sleep pressure in the third row was not fitted at all: the time constants are taken from one paper, the construction of the switch from another, and the agreement came out by itself. This is the weightiest of the checks.

% =====================================================================
\section{Slow defence: what was added and what it did}\label{sec:slow}
\figp[0.88\textwidth]{fig_slow.png}{\label{fig:slow}Four results of version~4. \textbf{a} --- the two stages of the defence: the induction signal is fast, the protein slow, and half the level is reached only by the third day. \textbf{b} --- the law of episode duration: it turned out to be a straight line rather than a saturating curve, and a test point outside the fit confirmed this. \textbf{c} --- the locking staircase: the rotation number stands at rational values and jumps; the inset shows the alternation pattern of cycle lengths on the fractional step 5:4. \textbf{d} --- the number of cells enters the model as an effective spread of resilience.}

The conservation law for the gaps (section~\ref{sec:law}) left exactly two ways out, and the second of them is variable defence: if after a load the cell strengthens its repair, the lower boundary of the window moves upwards by itself and the return does not require a large excursion of the noise. It is this way out that we built up.

\textbf{What exactly was added.} The peroxiredoxin repair enzyme (sulfiredoxin) is not up-regulated instantaneously: first an induction signal accumulates, and then the protein is synthesised from it. These are two first-order stages in series:
\[\dot m=\frac{1}{\tau_m}\Big(\beta\frac{u}{K_B+u}-m\Big),\qquad \dot s=\frac{1}{\tau_s}\,(m-s),\]
where $u$ is the peroxide, $m$ the signal and $s$ the accumulated protein, which adds to the constant defence. The steady state is the same as before: the window, its boundaries and the whole conservation law for the gaps are unchanged. Only the time course changes --- and it is precisely the time course that decides everything.

\textbf{Where the constants come from.} The time of the first stage is taken from the mRNA: 8.1--9.6\,h in three independent treatments.\tS\ That of the second from the protein, from the raw data of the same laboratory: a peak at hour 18, while the decay constant has never been measured in the published work, and comes out at anything from 24 to 88\,h depending on the method of processing.\tS\ The working value of 63\,h is the middle of that range, and it is our assumption, not a measurement.\tD\ The induction strength is drawn across sets: the measured fold change of protein accumulation runs $\times$2.1 / 2.7 / 4.2 / 7.0 at four oxidant doses.\tS

\textbf{What came of this --- and this is the principal negative result of the version.} The episode duration turned out to be a \emph{linear} function of the memory of the defence (fig.~\ref{fig:slow}b):
\[\text{duration}\;=\;1.28\;+\;0.97\,\tau_{\rm d}\quad\text{(days)},\]
with no saturation whatever. The functional form was not chosen by eye: three candidates were fitted to the points with memory below 100\,h, and the prediction was tested at the point 150\,h, which was not in the fit. The straight line beat both saturating forms, both on the residual of the fit and on this test.

Hence a consequence that closes an entire route. For the episode to fit within the clinical two days, a defence memory of about 12 hours is needed --- that is, \emph{less than that of the first stage}, and several times less than any of the measured values (24--88\,h). Passive defence does not miss the target by a hair's breadth: it is excluded by a factor of 2 to 10. Strengthening the repair does not in itself shorten the episode but lengthens it --- because the longer the accumulated protein persists, the longer the cell sits in an intermediate state instead of falling back quickly.

\begin{ineq}
Two first-order stages in series give the response to a step in the load
\[s(t)=1-\frac{\tau_s e^{-t/\tau_s}-\tau_m e^{-t/\tau_m}}{\tau_s-\tau_m},\]
that is, always with an inflection: there is no instantaneous response for any constants. For $\tau_m=12$\,h and $\tau_s=63$\,h half the level is reached by hour 57 (fig.~\ref{fig:slow}a). This is the reason why what must be accelerated is not the first stage but the \emph{gain}: the derivative of the response with respect to $\tau_m$ at the working point is small (on a single cell the acceleration gave $-1$ to $-5$\,\%), whereas that with respect to the amplitude is large ($-52$\,\%). The prediction was written down before the computation and was confirmed.
\end{ineq}

\begin{note}
From the fact that the defence lengthens the episode in the model it does \emph{not} follow that sulfiredoxin is harmful. One thing follows: a defence that only repairs what has been damaged does not explain the picture of the disease. It can be explained only by a defence that acts \emph{earlier and more strongly} --- that is, by priming (section~\ref{sec:prime}).
\end{note}

% =====================================================================
\section{Priming: the only route that reached the target}\label{sec:prime}
Priming is memory of a past exposure: a cell that has already met an oxidant responds to the next exposure more strongly and more quickly. It is inserted into the model in one line: the same signal that drives induction accumulates a trace $p$, and the induction strength is multiplied by $(1+A\,p)$. The ceiling $A$ is taken from measurement --- threefold.

\facts
\begin{itemize}
\item On a repeated exposure 24 hours later the response of the regulator is lower while the accumulation of the enzyme itself is \emph{higher} --- up to threefold --- and the time to half-maximum falls from about 8 to about 4 hours. The same laboratory and the same system that gave us the raw kinetic data.\tS
\item The effect depends on the master regulator of the oxidative response and does not depend on its auxiliary partner: that is, it is an amplification of the signal and not a side effect.\tS
\item Independently, in another tissue: accumulation of a protective enzyme under repeated exposures was followed over four days and was gene-selective.\tP
\end{itemize}

\textbf{What came out.} For the first time in the project the episode duration entered the clinical target: median 1.46 days against an upper bound of 1.665 (a one-sided estimate from order statistics), fraction of short episodes 0.363, paired test against the control without priming $p=5.2\cdot10^{-6}$. All four parts of the rule written down before the computation are satisfied --- and satisfied again on a sample of 2000 sets, six times larger.

It matters what this is \emph{not}. Priming was not fitted: its amplitude is taken as the ceiling from measurement, its time constant was scanned, and the acceleration of the first stage --- the second half of the measured effect --- was checked separately and gave almost nothing, exactly as predicted before the computation. One route remained, and it passed.

% =====================================================================
\section{The precursor: what produces it}\label{sec:prodrome}
Some days before a collapse, waking lengthens --- in the model this is visible as a rise of the curve ``length of waking against time to collapse'', normalised to each set's own healthy level. The question of what produces it was settled by paired comparison: the same sets, one sid, the difference in a single link only.

\begin{itemize}
\item \textbf{Slow noise in the load is not it.} When it is switched on the height of the precursor does not grow (paired sign test: a rise in 131 sets of the sample, $p=0.97$). The noise had been introduced on the justification that ``the model contains no quantity living on the scale of days''; that ceased to be true as soon as the two-stage defence appeared. The noise has been removed from the working set.
\item \textbf{The memory of the defence is partly it.} When it is taken away the height falls in 107 of the sets checked, $p\le0.05$. But only \emph{a fifth} of the excess goes away.
\end{itemize}

So there is a third slow quantity. The search narrowed to a single candidate by a simple enumeration of all the times of the model: the return of the peroxide and the relaxation of the load give 0.25--0.5 and 0.03--0.05 days; the return of $x$ is of the order of 0.01 days (and this is confirmed by direct measurement: in 73\,\% of sets $x$ returns within a day); the first stage of the defence is 0.5 days; the second is 2.6 days and has already been excluded as the sole cause. \emph{There are no other times in the equations.}

Hence the hypothesis recorded as rule F21: the third slow quantity is not a time constant but \textbf{accumulation}. A population of $J$ cells is filled with damaged ones one at a time, and the characteristic filling time is set not by a constant but by the \emph{rarity} of the transition, that is, by the distance of the load from the barrier in sigmas of the noise. Such a quantity has no $\tau$ of its own --- which is why a scan of constants could not have found it.

% =====================================================================
\section{What an episode consists of}\label{sec:tail}
All the other rules measure the duration \emph{between} sets. But for the choice of a mechanism what matters more is what the duration is made up of \emph{within} a set. An episode divides into two parts: the time during which the load stays above the top of the window, and the \textbf{tail} --- the time from the last such moment to the end of the episode, when the load is already normal but the cells are still damaged.

This distinction decides what the defence can get hold of at all. If the tail is large, repair is limiting, and it is repair that must be accelerated. If there is almost no tail, the episode lasts exactly as long as the high load persists, and repair is not limiting at all.

On the earlier version of the model (still without the two-stage defence) the tail occupied a quarter of the episode, and the association of duration with the time of high load was 0.94. A preliminary check on the \emph{present} model gives a sharper picture: association 0.87, tail about 7\,\% of the episode.\tD\ If this is confirmed on the full sample, the conclusion will be hard: \textbf{priming shortens the episode not by repairing faster but by raising the top of the window}, and thereby shortening the time of high load itself. This is a testable prediction, and it was written down before the computation.

\begin{note}
The earlier measurement of the tail was made on a model without the two-stage defence, and it cannot be carried over to the present one: it is precisely the defence that moves the boundaries of the window. The rule has therefore been rewritten, and the episode recorder is \emph{compared bit for bit} with the main computation on the same rows and the same sid: if a copy of the step diverges in even one episode, the rule declares itself unverifiable and prints no numbers.
\end{note}

% =====================================================================
\section{How many cells are in the population --- and why this turned out to be a parameter}\label{sec:J}
The number of cells $J$ was until recently chosen for computational convenience: 32 in some runs, 64 in others. It was assumed to be merely a resolution. It turned out to be a parameter, and it entered by three different routes.

\begin{enumerate}
\item \textbf{The tail of the resilience distribution.} The multipliers were taken as equiprobable quantiles $p_j=(j+1/2)/J$, so the extreme quantile drifts as $J$ grows: the more cells, the more vulnerable the most vulnerable one. \emph{Closed}: the probability grid is pinned at the ends, and at $J=32$ it coincides with the former one to the last digit, so the old runs remain valid.
\item \textbf{Quantisation of the episode threshold.} An episode was defined as ``the fraction damaged above $1/2$'', and the fraction takes only the values $j/J$, so the actual threshold is $0.625$ at $J=8$ and $0.516$ at $J=64$. \emph{Circumvented}: verdicts are issued on a continuous variable whose threshold is not quantised.
\item \textbf{The realised spread.} This one was found only on a sample of 2000 sets, when the measurement noise fell below the effect. The ends of the grid are pinned and indeed barely drift (the smallest multiplier is 0.8002 at $J=8$ and 0.8029 at $J=512$), but the \emph{variance} of a finite grid is smaller than the one specified and approaches it only in the limit: $\mathrm{std}(\ln v)=0.1226\,/\,0.1060\,/\,0.0980\,/\,0.0943$ at $J=8/16/32/64$, and $0.0910$ in the limit (fig.~\ref{fig:slow}d). From $J=16$ to $J=64$ this is 11\,\% in the effective sigma of the spread, and the model is extremely sensitive to it.
\end{enumerate}

The third route explains the dependence found quantitatively and indicates the correction: stretch the grid so that the realised spread equals the one specified for any $J$. The decisive check is written down before the computation: a block at $J=32$ with a substituted sigma is obliged to reproduce a block at $J=16$ with the ordinary one. If it does, the route is explained in full.

\begin{ineq}
A general lesson, not about our model. Discretisation of a continuous distribution by $J$ points can \emph{exactly} preserve the mean (by normalisation) and \emph{exactly} pin the ends (by the choice of the extreme quantiles) and still fail to preserve the variance: for a grid uniform in probability, $\mathrm{Var}\to\sigma^2$ only as $O(1/J)$ from below. If a model is sensitive to the variance of the spread --- and here the logarithmic sensitivity is about 2.5 --- then the number of grid points becomes a parameter of the model, and runs with different $J$ cannot be compared: that is a comparison of two different models.
\end{ineq}

% =====================================================================
\section{How the rhythm breaks}\label{sec:bifur}
A separate question, posed independently of the redox part: what happens to the alternation of sleep and waking as damage inhibits the sleep nucleus ever more strongly. The expectation was this: on approach to a critical value chaos should appear --- that is, the circadian rhythm breaks \emph{while} the 24-hour external pacemaker is preserved. This is cheap to compute: the switch is deterministic, and neither the cells nor the redox part are needed at all.

The right observable here is the \textbf{rotation number}: how many days of the external forcing fall on one cycle of sleep and waking. The rhythm does not break smoothly: the rotation number stands at rational values on finite intervals of the parameter and jumps between them (fig.~\ref{fig:slow}c):
\[1\!:\!1\;\to\;5\!:\!4\;\to\;3\!:\!2\;\to\;2\!:\!1\;\to\;3\!:\!1\;\to\;4\!:\!1\;\to\;\dots\;\to\;9\!:\!1\;\to\;\text{no sleep}.\]
The transitions fall at $kD\approx4.77$; 4.78; 4.79; 5.09; 5.16 and then ever more densely: five steps (5:1, 6:1, 7:1, 8:1, 9:1) fit between 5.172 and 5.178\,mV, after which sleep disappears altogether.

A separate run with a fine scan and refinement of the boundaries (fig.~\ref{fig:stairs}) shows that there are far more steps: \textbf{52 over an interval of 0.9\,mV, and in 34 of them the cycle lengths alternate}. Between the integer lockings stands a dense accumulation of fractional ones --- 5:4, 6:5, 4:3, 13:10, 11:8, 9:7, 7:5, 10:7, 16:11, 8:5, 12:7 --- and between those, points where the pattern does not repeat at all. The step widths fall with the denominator by many orders of magnitude: for 2:1 it is 0.3\,mV, for 3:2 about $8\cdot10^{-3}$, for 11:8 already $5\cdot10^{-6}$. A uniform scan cannot catch this: only the boundaries can be refined.

\figp{fig_stairs.png}{\label{fig:stairs}The devil's staircase in full, from a separate run with refinement of the boundaries (about 60 cycles per point after the transient was discarded). \textbf{Above} --- the whole interval: the flat steps of the integer lockings 1:1, 2:1, 3:1 and on to 9:1 packed together just before sleep disappears. \textbf{Below} --- a close-up of the region with the most steps: in orange, lockings with \emph{alternating} cycle lengths; in grey, regimes with no repeating pattern. The inset shows the pattern of the 3:2 locking itself: a long cycle of 47\,h, a short one of 25\,h, and so round again.}

\textbf{This is a locking staircase, not a period-doubling cascade.} In a doubling cascade the period is multiplied by two and takes no other values, and the Feigenbaum constant describes how the steps converge. Here the period adds one missed night at a time, and fractional steps stand between the integer ones --- these are Arnold tongues, the ordinary picture of an oscillation entrained by an external force. To invoke the Feigenbaum constant here would be an error.

\textbf{The most interesting thing is the fractional steps.} In them the cycle lengths are not equal: they follow a strict pattern and repeat. At $kD=4.772$\,mV the pattern has four cycles over five days: 25, 24, 24, 47\,h --- three nearly ordinary days, then one doubled, and again (the inset in fig.~\ref{fig:slow}c). At 4.779 the pattern has two: 47 and 25\,h. At 4.783, three: 46, 26, 24\,h. Checked for stability: the pattern is the same at 120 and at 400 days of computation, so it is not a transient.

This matters for the following reason. Such behaviour is not ``no rhythm'' but \emph{another rhythm}, with a period of its own of several days, and this is exactly how part of the clinical picture of disturbed sleep looks. The model produces it by itself, without any noise and without a single added equation --- simply as a consequence of entrainment. A prediction written down before the computation: with noise added, these windows should not disappear but blur into an \emph{intermittent} alternation --- the pattern holds, then breaks down, then holds again.

\textbf{The principal limitation is honest and unpleasant --- but it is not the one that stood here in version~4.} The ceiling is set not by the damage but by the coupling itself: identically $kD\le k$, so at $k=2$\,mV the step at 4.77\,mV is unattainable \emph{for any} damage. A coupling of no less than 4.77\,mV is needed, that is, 2.4 times stronger than the working one. At $k=5$\,mV the model \emph{does} reach the staircase: in sets that have fallen into the irreversible state, $kD=4.93$ was measured at a damaged fraction of 0.999 and waking of 0.974. So the limitation is not in the mechanism but in the calibration of a single constant.

\textbf{The previous edition named two ways out; the second is struck off.} ``Chronic damage holding the fraction of broken cells high'' does not work, and this is evident without any computation: in the broken sets the fraction is already at its ceiling, and the ceiling equals $k$. One way out remains --- the strength of the ``damage--sleep'' coupling. The earlier numbers in this paragraph (``twofold'', ``fourfold'') were wrong and did not even agree with each other; their root is analysed in section~\ref{sec:prec} as the seventh statistical trap.

\begin{note}
The check written down before the computation asked the wrong thing: it looked among the periods for the values 1, 2 and 4 --- and found them, but would have found them in a locking staircase as well. The verdict is left as it stands, and beside it are recorded the two parts that were missing: is there locking at a period that is not a power of two, and are there regimes with a pattern longer than one cycle? Separately, the scan grid had to be corrected: the 5:4 and 3:2 steps have a width of about 0.005\,mV and are invisible on a coarse grid --- the first run saw 1, 2, 4 and not a single fractional one.
\end{note}

\begin{ineq}
Why a staircase and not a cascade was what had to come out here. The switch is a relaxation oscillator with an intrinsic period of about 25\,h; the daily term is an external force of period 24\,h. For such a pair the ratio of the periods is the rotation number, and it locks at rational $n\!:\!m$ on intervals of finite width (Arnold tongues), varying between them along a devil's staircase. The width of a tongue falls rapidly with the denominator $m$: for 2:1 it is about 0.3\,mV, for 3:2 and 5:4 about 0.005, that is, sixty times narrower. A staircase can be distinguished from a doubling cascade by exactly one sign: the presence of locking at a period that is not a power of two.
\end{ineq}

% =====================================================================
\section{Asymmetric defence}\label{sec:asym}
\figp[0.88\textwidth]{fig_v5.png}{\label{fig:v5}Four results of version~5. \textbf{a} --- asymmetric defence: the measured level lies on the derivation $w(r)$, the ordering agreed exactly. The range is anchored by the two extreme points (circled); the three middle ones are predicted. \textbf{b} --- burstiness of the load: with the mean and the variance pinned, the collapse rate falls by 40\,\%, whereas the fraction of cells that were damaged saturates and sees nothing; the choice of the measured quantity decided the outcome of the rule. \textbf{c} --- the width of the rescued band against the damaged fraction: two scans give opposite signs, so the band is not a function of $D$. \textbf{d} --- the cost per row doubles for the one-operation-at-a-time computation once the working set no longer fits in cache, and does not change at all for the fused kernel.}

The second stage of the defence had all along been written with a single constant: $\dot s=(m-s)/\tau_s$, 63 hours both on the rise and on the decay. Biologically these are two different processes --- induction of synthesis and degradation of protein --- and their equality was an assumption we had simply never tested.\tD\ Version~5 separates them: $\tau_{\uparrow}$ on the rise, $\tau_{\downarrow}$ on the decay, with the ratio $r=\tau_{\downarrow}/\tau_{\uparrow}$.

\textbf{A derivation yielding a prediction before any computation.} Under a daily cycle the input $m$ is two-level to good accuracy: the high level holds for a fraction $f$ of the day, the low level for the rest. For a leaky integrator with different constants, when both are much larger than a day, the steady mean stands between the levels at the point
\[w(r)=\frac{f\,r}{f\,r+1-f},\]
that is, $s_{\rm mean}=m_{\rm lo}+w(r)\,(m_{\rm hi}-m_{\rm lo})$. The asymmetry turns the integrator into a \emph{peak detector} for $r>1$ and into a \emph{trough detector} for $r<1$. The fraction $f$ is taken from the run (it equals the measured fraction of waking, 0.705) rather than assigned, so the prediction contains not a single fitted number.

\textbf{What came out.} At $r=0.33/0.5/1/2/3$ the derivation gives $w=0.441/0.544/0.705/0.827/0.878$, while the measured level of the defence is $0.0575/0.0644/0.0792/0.0958/0.1065$. The ordering agreed \emph{exactly}: rank association $+1$ at an exact permutation $p=0.0083$ (fig.~\ref{fig:v5}a). The consequence for collapses was confirmed as well: the fraction of cells that were damaged falls from 0.486 to 0.457, the trend test gives $z=1.67$, $p=0.048$, and the episode duration shortens from 7.74 to 6.88 days.

\textbf{The conclusion that changes the status of the earlier model.} The equality of the rise and decay constants was not a neutral simplification but a \emph{choice of regime}. Different $r$ give a different mean defence under one and the same instantaneous equations, and the direction of the effect is predictable from two lines of algebra. So either the equality must be justified by measurement, or $r$ must remain a parameter of the model on the same footing as $\tau_s$.

\begin{note}
The quantitative part of the rule remained unverifiable, and the reason lies in the \emph{construction of the rule itself}, not in a shortage of sets. The equivalence tolerance $3\,\mathrm{SE}=0.0241$ and half the range of the curve, $0.0245$, almost coincided, so the gap between the threshold and the alternative is 0.0004 at an MDE of 0.02. No number of sets cures this: 29 times more would be needed. What cures it is leverage --- widening the scan of $r$ from 0.1 to 10, which increases the range of the curve 1.7-fold. This is the price recorded honestly: we know \emph{what} we did not verify and \emph{how much} it costs.
\end{note}

% =====================================================================
\section{Three phases: the rescued band}\label{sec:band}
The cells of one set differ in exactly \emph{one} scalar --- the load multiplier $v$; everything else is common to them. So the state of a cell is a one-parameter family $x(t;v)$ under a common load history. Were this family monotone in $v$, the damaged cells would always be a clean tail $\{v\ge v^*\}$, the damaged fraction would be computed in closed form as $1-F_v(v^*)$, and the per-cell block would not be needed at all.

\textbf{Monotonicity does not hold, and what breaks it is the slow defence.} On a fine grid of 600 nodes in $v$ under three load profiles, including a deliberately adversarial one, the measurements are:
\begin{itemize}[leftmargin=1.1em,itemsep=1pt,topsep=2pt]
\item with the second stage \emph{frozen}, monotonicity is exact: zero violations, the largest exactly $0.000$;
\item with a live defence the ordering breaks on 11.9\,\% of the steps, the largest violation being 0.973 --- that is, a complete reversal: the more vulnerable cell is intact while the less vulnerable one has collapsed;
\item priming is not the cause but an amplifier: at $A=0$ the violation remains but falls from 0.973 to 0.0184;
\item the ordering is preserved in no coordinate: $s$ --- 27.8\,\% of the steps, $m$ --- 32.5\,\%, the priming trace --- 32.7\,\%, the total defence --- 100\,\%.
\end{itemize}

\textbf{The mechanism.} The more vulnerable cell is oxidised earlier, is primed more strongly, accumulates more protective protein --- and by the time a large excursion arrives it has \emph{more} defence than its less vulnerable neighbour. This is hormesis at the level of the population: part of the cells survive precisely because they suffered earlier.

\textbf{There turned out to be three phases.} The boundary layer is not smeared but connected: in a pilot its width held at 0.085 of the population as a structure. In the production computation at $J=32$ the width of the band (the fraction of intact cells lying above the first damaged one) equals $0.00732\pm0.00042$ --- seventeen sigmas from zero. The cause was named correctly and verified: with the defence \emph{frozen} the band becomes exactly \emph{zero}; with priming switched on it is 1.82 times wider ($0.01331$); and a short memory of the second stage does not change it at all ($0.00712$ against $0.00732$).

\begin{ineq}
The shape of the band was predicted by geometry: the boundary layer must follow $\varphi(\Phi^{-1}(1-D))$, that is, be widest at $D=1/2$. \textbf{The prediction is refuted, twice over and with different signs.} Over the scan of the sleep coupling the rank association is exactly $-1$: the band \emph{narrows} as the damaged fraction grows, from $0.0135$ at $D=0.10$ to $0.0021$ at $D=0.34$. Over an independent scan of the daily range, which was not opened for this question, it is exactly $+1$: the band \emph{widens}, from $0.0095$ at $D=0.064$ to $0.0191$ at $D=0.162$. The ranges overlap, and at one and the same damaged fraction the width differs twofold depending on \emph{which knob} was used to change that fraction (fig.~\ref{fig:v5}c). The claim about the shape is withdrawn in full, as had been recorded in advance.
\end{ineq}

\textbf{What this means for the theory.} A closed formula for the damaged fraction in terms of a single boundary is illegitimate with a live defence. The per-cell block may be simplified only with an \emph{explicit} boundary layer --- three configurations, not two. And the width of that layer is not determined by the observable damaged fraction: it is set by something third, most probably the time the population has spent right at the tipping point.

% =====================================================================
\section{Mean field: what follows from it and what will not}\label{sec:mf}
The structure of the model looks like an invitation to mean field: all-to-all coupling through two scalars (the damaged fraction and the mean damage) which enter only the sleep switch. An analysis by method gives the following picture.

\textbf{Hartree, ordinary mean field.} Does not work, and it is now known why: the damaged fraction is not a function of the mean but a functional of the distribution, and precisely at the threshold, where the distribution is bimodal because of the bistability of the redox switch. There is nothing to average there.

\textbf{Dynamical mean field (DMFT).} Formally \emph{exact}: all-to-all coupling through a scalar is infinite dimensionality, and in the limit $J\to\infty$
\[D(t)=\mathrm{P}\big(v:\;x(t;v)<x_{\text{thr}}\big)\]
with no approximation whatever. The limit exists and has been measured: the corrections from the realised spread of the grid converge as $1/J$, and extrapolation from the last two points gives 0.1101. But DMFT by itself yields no economy: it reduces the problem to a single-cell one plus a quadrature over $v$, and the whole cost moves into the number of quadrature nodes.

\textbf{Level set.} Would give a complete solution --- one boundary instead of $J$ cells. Dead with a live defence, \emph{exactly correct} with a frozen one (section~\ref{sec:band}).

\textbf{Born--Oppenheimer.} The separation of scales is real: 63 hours against the fast relaxation. On the fast scale the damaged fraction is closed and the slow variable moves the boundary --- but it moves it \emph{unequally across cells}, and hence the band.

\textbf{Multiconfiguration approximation.} The route that is needed, and the number of configurations is now known from measurement: three, plus two moving boundaries.

\begin{note}
\textbf{Decision on application: the reduction is not to be used as a speed-up.} The gain is estimated at three- to fourfold (32 nodes against 8--12) against a run that already fits in three hours. The price is an edit to the model which cannot be checked by identity, which voids the base of verdicts and which requires a convergence study in $J$ of its own. And, most importantly, the reduction would coarsen precisely the boundary layer, that is, the one new phenomenon. What pays off at once and for free is the closed level-set formula with a frozen defence as a source of \emph{analytic expectations}: the model is not touched, the cache is not devalued, and the thresholds acquire a derivation from expectations and variances.
\end{note}

% =====================================================================
\section{Burstiness of the load: a hypothesis that did not pass}\label{sec:burst}
Two discrepancies with the arithmetic of the hazard pointed in \emph{opposite} directions, and that looked like a lead. The observed increase of the collapse rate due to surges turned out to be tenfold \emph{stronger} than predicted; the observed fraction of cells damaged at a doubled daily range, threefold \emph{weaker}. Neither the convexity of the hazard nor the replacement of the density by the exceedance probability can produce such a thing: both corrections move the prediction in one direction. What remained was a mechanism absent from the hazard altogether --- lag of the defence: a short surge arrives at an unadapted cell, whereas a range raised for a day is partly compensated by the defence in time.

\textbf{How this was tested without a single fitted number.} To write an adaptation factor into the predictor would mean introducing a parameter derivable only from the algebra of the fold, that is, in effect fitting it. The question was therefore posed as a comparison in which everything extraneous is pinned: the strength of the slow noise sets the \emph{variance} of the load at an unchanged mean (the mean of the lognormal is compensated within the model itself), and the correlation time sets the \emph{time scale} at an unchanged mean and variance. A scan over the correlation time isolates exactly one thing: whether the defence has time to adapt.

\textbf{The rule as written is not satisfied.} The fraction of cells damaged, over ordered correlation times, gave $z=0.61$, $p=0.27$: there is no trend.

\textbf{And the defect here is mine, not the model's.} The fraction damaged --- ``was the set damaged even once in 60 days'' --- \emph{saturates}: it runs $0.599\to0.623\to0.626\to0.627\to0.614\to0.588$, that is, as a hump, and has almost no sensitivity. The collapse rate on the same blocks, meanwhile, falls monotonically from 0.0550 to 0.0328 --- by 40\,\% --- and the episode duration grows from 5.37 to 8.04 days (fig.~\ref{fig:v5}b). An estimate from the published table gives quantitative agreement as well: the share of the fall falling between the two constants of the defence equals 0.523 against the 0.424 expected under a fall uniform in the logarithm.

\begin{note}
This estimate \textbf{is not a verdict} and has not been entered in the ledger. It was computed after the data had been seen, so its threshold is fitted to them. For F24 to receive an honest verdict, the rule is rewritten in terms of the collapse rate and rerun on a \emph{different} configuration --- a different sid and a different strength of the slow noise --- with thresholds written down in advance. The verdict ``not satisfied'' of 17.10 remains in the ledger as it stands; the defect is noted beside it. This is the same discipline by which we rewrite nothing after the fact, our own misses included.
\end{note}

\textbf{The lesson is wider than the rule itself.} To the four statistical traps of version~4 a fifth was added, and it is about the choice of the \emph{measured quantity}, not of the criterion: a saturating quantity quenches a real effect, and no statistics will correct that. The mechanism acts on the rate of events --- so the rate must be measured, not the fraction of those to whom the event happened at least once.

% =====================================================================
\section{The dopamine loop: a latch, not a knob}\label{sec:dopa}
Dopamine appeared to have been cut out of the model, while the results remained reasonable. The premise turns out to be wrong: dopamine is not cut out but \emph{bracketed at both ends}. Its contribution to the oxidative load sits inside the constant influx (catabolism of dopamine yields peroxide); its action on the balance of excitation and inhibition is carried by the activity itself through the ratio $E/E_r$; and Grace's loop takes as its input exactly what the model has at its output. A feedback loop was in the model already --- carried by glutamate: damage $\to$ disinhibition ($g=1-0.4\,D$) $\to$ a rise of $E/E_r$ $\to$ a rise of the load $\to$ damage.

The one argument for adding it was not ``it exists in the brain'' but a measured failure: the working region does not reach the locking staircase. The pilot showed that this argument too is wrong --- but the loop turned out to be substantive for a different reason.

\textbf{What was added.} One scalar, a tone $d$, driven by the damaged fraction itself (Grace's loop runs from the damage through the hippocampus and the VTA, bypassing local cortical activity), with a constant $\tau_d$, and back into the load as a factor $(1+a\,d)$. At $a=0$ the loop is open and the model is reproduced \emph{bitwise}.

\facts
\begin{itemize}
\item The direction is confirmed: catabolism of dopamine yields peroxide; in Gclm-knockout mice, blockade of dopamine reuptake during postnatal development --- but not in adulthood --- produced a persistent fall in PV-cell density. The vicious circle is described as a ``dynamic feed-forward loop'' in a review in Molecular Psychiatry.\tP
\item But there is \emph{not one} quantitative value: neither the rate of ROS formation from dopamine, nor the time constant of the loop. The effects are described in the words ``permanent'', ``long-term'', ``persistent'', and are tied to the developmental window.\tS
\end{itemize}

\textbf{Verdicts (2400 sets).} (a) satisfied: the fraction of irreversibly broken sets is 0.7076 with the loop against 0.0829 without, at an MDE of 0.037. (b) \textbf{not satisfied}: $-0.2924$ --- the loop breaks \emph{fewer} sets than its static equivalent (factor 8, the loads matched to 87\,\%), which breaks 1.000 against the loop's 0.708. (c) satisfied: under the two-phase protocol the loop releases 0.0024 of the broken sets against 0.9936 for the open model. (d) \textbf{not satisfied} --- the price, and it is heavy.

\textbf{The substantive conclusion: the loop is neither a knob nor an amplifier but a \emph{latch}.}
\begin{itemize}
\item \emph{Not a knob}: a knob gives no hysteresis. The open model, when the load is removed, releases \emph{all} the broken sets ($0.9936\to0.0000$ exactly); the loop at $a=10$ releases 0.24\,\%.
\item \emph{Not an amplifier}: it breaks \emph{less} than its static equivalent, because it is selective --- its factor is high only in cells already damaged and stays near unity in healthy ones.
\item \emph{A latch with nucleation}: at a load factor of 0.6 the open model breaks \emph{exactly zero} sets, whereas the loop breaks 2.3\,\% at $a=3$ and 3.7\,\% at $a=10$, having raised the load from 0.161 to 0.285 by itself. A random dip latches permanently where the static load does not suffice at all.
\item \emph{Dose dependence}: at $a=3$ it releases 3.8\,\%, at $a=10$ only 0.24\,\%. So below there is a critical strength at which the latch does not yet engage.
\end{itemize}

\textbf{Is there a window where hysteresis is already present and episodes are still present?} That is the clinical picture --- a chronic disease with exacerbations --- and the question was posed as a separate rule, valuable in either direction. The answer is \textbf{no}. The latch engages at a loop strength of 1.11 and the episodes die out at 0.18: a sixfold gap, and a structural one. Monotonicity of both quantities in the strength is confirmed ($0.9958$), so there is nothing between them to look for; the control at strength zero gives retention of exactly 0.0; and under a softer breaking protocol the latch merely moves to 0.98 without opening a window.

\begin{ineq}
The cause is named and it is not a matter of parameters. The load is a quantity \emph{per set} --- one for all the cells --- so a loop through it moves the whole population at once, and what the population loses is precisely its heterogeneity, the only thing that gave episodes somewhere to live (section~\ref{sec:pop}). A loop of this kind therefore describes not the disease but its irreversible outcome. The time constant $\tau_d=1$ day is not taken from the literature, and the rule was built so that the answer does not depend on it: it asks about the \emph{difference} from the static equivalent, not about the magnitude of the effect.
\end{ineq}

% =====================================================================
\section{The capacity loop: episodes rarer but longer}\label{sec:cap}
The load loop failed structurally, and the cause pointed at the repair: a loop must act \emph{per cell}, so that a damaged cell loses something its neighbour does not. The capacity of the defence is such a quantity. While a cell is damaged its capacity $N_0$ is drawn towards $N_0(1-a_n)$ with a fast constant, and on recovery returns to $N_0$ with a slow one. The damaged cell loses capacity, the neighbouring one does not --- the population stays heterogeneous, and the episodes have somewhere to live.

\facts
\begin{itemize}
\item This part is measured: oxidative damage to PV cells lowers parvalbumin itself, tears the perineuronal net and strikes the mitochondria.\tS
\end{itemize}

The depth to which the load is removed stands as an \emph{axis of the scan and not as a constant}: in the pilot the answer about hysteresis moved from 0.037 to 0.648 depending on how deeply the load was lifted. A constant would have given an answer about the measurement regime rather than about the model.

\textbf{The falsifier broke, and taking it apart gave more than the pilot itself.} Four defects were found, two of them in our computation and two in the construction of the rule.

\begin{enumerate}
\item \emph{The kernel diverged from the reference, and not one pin touched it.} The fused kernel aimed at the raw $N_0$ instead of the reduced capacity $N_0\cdot n_0$. The reference gave exactly $0.85\,N_0$, the kernel $0.88886\,N_0$: the capacity crept back towards the unreduced value with the fast constant. The defect appeared \emph{only} for $n_0<1$, and all the pins ran at $n_0=1$, where the two agree identically. Two pins were added --- a static capacity, and a static capacity together with the loop (they multiply inside the relaxation target, and that is exactly where the error was invisible) --- and it was verified that the new pins \textbf{fail} without the correction: 52 fields of 67 diverge, the largest discrepancy being $3.2\cdot10^{33}$. A pin that passes before the fix is worth nothing.
\item \emph{The cache outlived the kernel.} The kernel source was not part of the code-version hash, so an edit to the kernel did \emph{not} devalue the block cache. Corrected, at a stated and one-off price: 96 cached blocks devalued, one full recomputation needed.
\item \emph{Two thresholds were unattainable by construction.} Two parts required $p\le0.05$ from an exact permutation test over \emph{three} scan points. The smallest one-sided $p$ attainable from that test is $1/n!$, which at $n=3$ is $0.167$: the threshold could not have been met by any behaviour of the model. The price is visible in one of them: $\rho$ came out at exactly $+1.000$ --- perfect monotonicity, the model behaving exactly as predicted --- and the part was recorded as ``not satisfied''. That is a false accusation against the model for failing a check that cannot be passed. Corrected two ways at once: the grids were widened to five points (smallest attainable $p=0.0083$, the span of the scan unchanged), and a check of attainability was built into the rule --- if $1/n!>0.05$ the part is honestly declared \emph{unverifiable} rather than failed.
\item \emph{The comparison point was taken as the ``nearest'' element of a grid that did not bracket it.} The falsifier compares the loop against a static reduction of the capacity \emph{of the same mean magnitude}. The loop reached a capacity of 0.694 of the baseline while the grid ended at 0.85: the rule silently took the extreme point (0.632 against 0.444) and issued a verdict anyway. The answer obtained was about the depth of the reduction, not about the loop. Corrected: the grid was widened to 0.60--1.0 and now brackets with margin on both sides; the agreement tolerance was assigned \emph{before} the computation from the construction of the scan --- half the smallest grid step in capacity (wider, and ``agreement'' stops being agreement; narrower, and it is unattainable on any grid); and failure to bracket makes the part unverifiable rather than failed.
\end{enumerate}

\textbf{The rule rebuilt.} The replacement rule contains no test against zero at all: every part is a \emph{matched difference} --- the loop against a static equivalent levelled on the first half of the record, paired within a set, on the same sid. The verdicts: the machine control reproduces the analytic expectation $\mu\langle N_0\rangle$ to $1.9\cdot10^{-13}$; progression over the matched static is $+0.0231$ at an MDE of $0.0096$, and it replicates on three independent noise realisations ($+0.0223$ and $+0.0211$); retention under the softest release is 0.561 and grows with the softness monotonically; episodicity holds at 0.753 of the open model. One part did not pass: the course of the hazard over the matched static, $1.937$ --- it is recorded as not satisfied, and no cheap rescue for it exists.

\textbf{What the loop does.} It does not make collapses more frequent; it makes them longer. The duration goes $13.1\to25.6$ days while the rate falls $0.042\to0.033$ per set-day at risk, and the two move in opposite directions by construction of the measurement, not by fitting.

\begin{note}
The honest caveat: $25.6$ days lies far above the clinical two days, so the working point of the loop is overstated. What is comparable here is the \emph{direction and the ratio}, not the magnitude. This matters for how the prediction of section~\ref{sec:pred} is to be read: in longitudinal cohorts, as damage accumulates the episodes should become rarer and longer, not more frequent.
\end{note}

\begin{ineq}
The methodological lesson is wider than the rule. Identity of the fused kernel with the reference is a \emph{claim that is verified}, not a property that holds: the area of the verification equals the area of the pins, and a hole in it is invisible from inside --- it becomes visible only when a rule asks the model about something the pins never asked. And a cache must live exactly as long as the engine that filled it: a cache cannot verify identity, so it is not entitled to assume it.
\end{ineq}

% =====================================================================
\section{The sleep switch solved on paper}\label{sec:anal}
The whole of section~\ref{sec:bifur} was obtained by computation. The question then arose whether the same could be had on paper, and it can: with the fast pair of voltages eliminated adiabatically, the switch is a relay with hysteresis, and its return map admits a closed form up to two transcendental equations. The reduction settles, among other things, the question that the scan could not settle --- whether there are invariant tori in the window.

\textbf{The folds in closed form.} The steady state of the fast pair admits a \emph{parametric} solution: taking $V_m$ as the parameter,
\[Q(V_v)=\frac{A_m-V_m}{|\nu_{mv}|},\qquad D_v(V_m)=V_v-\nu_{vm}Q(V_m),\]
with $Q^{-1}(q)=\theta+\sigma_Q\ln\bigl(q/(Q_{\max}-q)\bigr)$, and the folds are the extrema of $D_v(V_m)$. Evaluating: $1.4503$ and $2.4635$\,mV against the $1.450$ and $2.463$ quoted by the computation, to four significant figures \emph{without integration}.

\textbf{The return map.} In the relay limit the somnogens obey linear equations with closed-form solutions, and switching occurs at two \emph{moving thresholds} $g_\pm(t)$. If $t_n$ is a sleep onset, then $t_{n+1}=f(t_n)$ is given implicitly by
\[g_+(t_n)\,e^{-(t_w-t_n)/\tau_s}=g_-(t_w),\qquad 1-\bigl(1-g_-(t_w)\bigr)e^{-(t_{n+1}-t_w)/\tau_w}=g_+(t_{n+1}),\]
both of the form ``exponential $=$ sinusoid''. The free-running period follows in closed form: $T_s=1.337$\,h, $T_w=6.683$\,h, $T_0=8.020$\,h against $8.878$\,h from direct integration. The 10\,\% discrepancy is not an error of derivation: switching does \emph{not} occur at the fold. Sleep onset is observed at $D_v=2.5026$ against the fold at $2.4635$, because the relaxation time diverges at a saddle-node and a slow passage completes after the bifurcation point. The overshoot was measured over two decades in the fast time scale and vanishes with exponent $0.64$ and $0.59$, against the $2/3$ expected for slow passage through a saddle-node. The closed form is therefore the exact relay limit, and the model stands 10\,\% away from it --- which must be stated rather than concealed.

\textbf{The tongue edge without a scan.} The $1{:}1$ orbit is a fixed point of $f$ shifted by one day, and the edge is its saddle-node: $f(t^*)=t^*+24$ together with $f'(t^*)=1$ gives three equations in three unknowns. The result is $4.7651$\,mV; the independent scan of the rotation number gave $4.77$\,mV.

\textbf{Arnold tongues do not apply at all.} The somnogens rise towards unity from below, so a sleep onset is possible only where $g_+(t)<1$. The fraction of the day admitting one falls from 0.789 at zero coupling to 0.166 at $4.80$\,mV and to 0.016 at $5.18$, where it vanishes --- and that is how sleep is lost. At $4.80$\,mV the domain is an arc of length $3.97$\,h and its image an arc of length $0.993$ inside it. A degree-one circle map has an image covering the whole circle; here the image occupies 4\,\% of it. The map is piecewise smooth with discontinuities: at $4.88$\,mV the jump is $22.47$\,h against a median step of $10^{-4}$\,h --- almost exactly one day, that is, a skipped night. \emph{The locking staircase is generated by branch switching, not by mode locking of a smooth circle map.}

\textbf{No invariant tori, and this is an inequality rather than an observation.} The right quantity is the Lyapunov exponent \emph{along the orbit}. Iterating the map and discarding the transient, over the whole window $\lambda$ lies between $-3.54$ and $-1.62$, with rotation numbers $1$, $3/2$, $2$, $3$ reproducing the staircase. Two conclusions follow: chaos is excluded, since it requires $\lambda>0$; and a smooth invariant torus is excluded, since a quasiperiodic attractor requires $\lambda=0$ exactly, whereas the distance to zero here is at least $1.6$. The attractor at every parameter examined is a strongly attracting periodic orbit. This is the third independent confirmation of the reduction, after the folds and the tongue edge.

\begin{note}
Reaching the right quantity required discarding two wrong ones. The first was the locked fraction on a grid: a locked fraction below unity was taken to imply that tori occur on a parameter set of positive measure. The inference is invalid --- the complement of the locked set is \emph{dense}, so a finite grid finds unlocked nodes whether the complement has positive measure or measure zero, and misses them in both cases with a probability depending only on the resolution. The criterion has no power in either direction, and the honest content of ``240 of 240 nodes locked'' is a lower bound on the measure of the locked set, $0.988$ at 95\,\% confidence, and nothing more. The second wrong quantity was the maximum of $|f'|$ over the domain: at tangential crossings the derivative has a square-root singularity, and a grid samples points the attractor never visits. The general lesson: on a Poincar\'e section, a quantity averaged \emph{along the orbit} and a quantity taken over the \emph{whole domain} answer different questions, and the latter need not concern the attractor at all.
\end{note}

\begin{ineq}
Measure in \emph{phase} space is a different question and does not arise here. All of the above concerns parameter space. The measure of the attractor itself --- a thick attractor of positive Lebesgue measure, whose dimension equals that of the system --- belongs to a harder circle of problems; Ilyashenko constructed examples for step skew products. The present system sits at the opposite pole, and this is measured rather than assumed: the fast time scales (10\,s) against the slow ones (hours) give a strongly negative divergence, and the two fast dimensions collapse at once --- the control established one-dimensionality of the return map to $4.5\cdot10^{-13}$, eight trajectories from different initial somnogen values falling onto a single curve to machine precision. The one-dimensional map then contracts onto a periodic orbit with $\lambda\le-1.6$, so the attractor on the section is a finite set of points and, in the extended phase space, a closed curve. This is structural: approaching thickness would require the divergence to cease being negative somewhere, that is, the fast time scales to become comparable to the slow ones --- but the entire reduction to a relay, and with it the folds and the circle map, rests on their separation. What formally remains open is an exceptional parameter carrying a Denjoy-type Cantor minimal attractor; the argument against is again structural rather than statistical, since such an attractor also requires the exponent at zero.
\end{ineq}

\textbf{What this buys in computation.} Constructing the map requires two scalar roots per cycle instead of 1600 integration steps per day, and the tongue edge requires solving three equations instead of a bisection. The scan took 1686\,s; the analytic route yields the same numbers in seconds, up to the stated 10\,\% offset. It applies in particular to the part the scan returned as untestable: the exponent along the orbit is obtained in twelve iterations of the map and answers its question directly, whereas a grid of 240 nodes could not answer it in principle, since it searched for aperiodic nodes where the attractor is periodic at every parameter.

% =====================================================================
\section{How we compute now}\label{sec:method}
A rule is fixed before the computation; this was so from the start. Version~4 added the second half of it: \emph{every criterion must have a known null value, a known noise level at our number of sets, and a known power}. Otherwise a verdict speaks about the sample and not about the model.

Formally: a part of a rule is deemed verifiable only if the distance from the null value to the threshold is no less than the minimum detectable effect
\[\mathrm{MDE}=(z_{0.95}+z_{0.80})\,\mathrm{SE}=2.4865\,\mathrm{SE}.\]
Otherwise the verdict is ``unverifiable'', and \textbf{symmetrically}: both where the rule would have failed and where it would have passed. Beside it is printed how many sets are needed for the threshold to become detectable. For claims of the form ``the observation is consistent with the prediction'' the tolerance is taken as 3\,SE: it must be wider than the MDE, otherwise such a claim is irrefutable in principle.

This is not a formality: four rules turned out to be unverifiable at any number of sets, because they were measured over eight cells, where the error of the rank association equals $1/\sqrt{7}=0.38$ while the threshold was set at $-0.5$. This is cured not by the sample but by subdividing the cells.

\textbf{Four statistical traps found along the way}, each of which cost a run of its own:
\begin{enumerate}
\item \emph{The median across sets of the fraction of rare events} is always 0 or 1: it collapses to zero as soon as fewer than half the sets have any events. Three rules fell on this, not on the model.
\item \emph{The fraction of short intervals} is identically $1+\mathrm{CV}^2$ over the sets' rates and sees no physics at all. Replacement: the $\mathrm{CV}^2$ of the intervals \emph{within} a set, where under a refractory period $\mathrm{CV}^2=(1-t_r/g)^2$ and the null is exactly one.
\item \emph{The direction of a sign test}: one and the same function applied to ``smaller is better'' and ``larger is better'' gave $p=0.99$ where the correct direction gave 0.038. The direction is now a mandatory argument, and the self-check runs both.
\item \emph{Zero error by construction}: the column ``in units of the set's own healthy level'' equals one for every set in the healthy cell, the spread is zero, and any verdict was declared significant.
\end{enumerate}

\textbf{Reproducibility.} The computation is distributed across files: the model, the statistics, the rules, the block assembler. Each rule \emph{declares} which blocks it needs; identical blocks of different rules are merged by hash and computed once (of 78 requested, 42 are computed). The result of a block is placed in a cache whose key includes the hash of the model's source code: an edit to the model devalues the cache by itself, and old blocks are not mixed into new ones. A separate test requires \emph{bitwise} agreement with the former monolith at pinned points. Verdicts are written to the ledger with the date and the block hashes and are not revised after the computation: the unfitness of a rule is noted beside it rather than corrected after the fact.

\begin{ineq}
A useful small point which was expensive here. The range of $n$ values of a quantity with relative error $r$ is itself a random variable: for three independent values $\mathrm{E}[\text{range}]=1.693\,r$. A threshold ``the spread over $J$ is no more than 10\,\%'' is therefore meaningful only for $r<0.059$. At 300 sets $r=12$\,\%, so the rule was unverifiable; at 2000 it became verifiable --- and immediately failed, revealing the third route by which $J$ enters the model.
\end{ineq}

% =====================================================================
\section{Accuracy of computation and verifiability of claims}\label{sec:prec}
Version~5 made the computation three times faster and at the same time established what in it may be believed.

\textbf{Where the speed-up came from.} The computation was limited not by arithmetic but by memory bandwidth: every elementary operation on the array of cells is a separate pass, and a step made about forty-five such passes. The whole per-cell block was rewritten as a single fused loop in which the intermediate quantities live in registers. The gain at the working size is \textbf{3.3-fold in wall time}: a full run of 123 blocks at 2400 sets takes 3 hours 05 minutes instead of nearly four hours previously at a smaller volume. The mechanism is visible in fig.~\ref{fig:v5}d: for the one-operation-at-a-time computation the cost per row doubles between 2000 and 5250 rows --- the working set ceases to fit in cache --- and holds thereafter, whereas for the fused kernel it is flat at all sizes, because its working set is a single row.

\textbf{What proves that this is the same model.} The fast path was not rewritten by hand but \emph{assembled from the source} of the reference by text substitution with a mandatory uniqueness check on every fragment: on any edit to the model the assembly fails rather than diverging silently. The comparison is made at two levels --- all 63 fields of the result at five points (including both sides of the asymmetry; the counting fields agree bitwise, the worst discrepancy of the continuous ones being $9.2\cdot10^{-15}$) --- and separately the whole set of rules, run twice with a diff of all verdicts.

\textbf{What to believe in the time step.} The quantities split into two classes. The continuous means converged with margin: at a step twice as coarse they hold within 0.5\,\%. The counting ones did not: the collapse rate falls monotonically under coarsening of the step in all conditions, and Richardson extrapolation gives a limit \emph{2\,\% above} the value at the working step. The cause is not the moment of threshold crossing: sub-step interpolation refines it but yields only 0.009\,\%. The bias comes from missed sub-step excursions and from the error of the trajectory itself. The conclusion for reading the results: thresholds on counting quantities finer than about 3\,\% may not be believed, and this is recorded as a known limit rather than painted over.

\textbf{The seventh trap, and it has already found its way into the text of this brochure once.} \emph{The mean over a bimodal population.} At the working point the sets fall into two heaps: either broken irreversibly (damaged fraction above 0.9) or untouched (below 0.01), with almost nobody in between. The mean over such a population is the \emph{fraction of broken sets}, not the state of the system, and it may not be read as a state. This is precisely how the claim arose in version~4 that the model does not reach the locking staircase: a mean over the mixture was being compared, whereas in the broken sets the damage stood at unity even then (section~\ref{sec:bifur}). The rule for reading: wherever a damaged fraction appears in the tables, the fractions of sets above 0.9 and below 0.01 must be kept beside it, or the mean will deceive.

\textbf{The fifth and sixth traps.} The fifth is the choice of the measured quantity (section~\ref{sec:burst}). The sixth is an \emph{irrefutable claim of agreement}: an equivalence tolerance exceeding the curve's own range will pass any data. A rule is now obliged to check this too, and on violation returns ``unverifiable'' with a statement of how many sets are needed.

\begin{ineq}
A separate device which paid off at once. Five parts of rules were unverifiable not for want of sets but by an \emph{information limit}: the rank association was measured over eight cells, its error is $1/\sqrt{n-1}=0.38$, so the detectable threshold is 0.94 against the $-0.50$ recorded. There is nothing to subdivide the cells into: the curve is indexed by wakings, and there is about one per day. The way out is to compute the slope for \emph{each set} over its own cells and to test the spread between sets, where $n$ is in the hundreds. The threshold must then be recomputed: the curve of a single set is noisier than the pooled one, and the association is attenuated by the classical factor
\[\rho_{\text{set}}=\rho_{\text{pooled}}\big/\sqrt{1+\lambda},\qquad \lambda=\frac{\text{variance of the cell noise}}{\text{variance of the signal}},\]
with $\lambda$ estimated from the same data \emph{before} we looked at the slopes. On synthetic data with a known answer the formula is accurate to 1--5\,\%. In the event $\lambda=33.8$, the threshold $-0.50$ was recomputed to $-0.085$, and the MDE fell from 0.94 to 0.076 --- twelvefold --- and four parts out of five closed on data already computed.
\end{ineq}

\textbf{Two rules were rebuilt along the way.} The rule on the composition of an episode lived on a \emph{second} integrator, a copy of the step inside the rule itself, with its own self-comparison against the reference. The comparison failed as soon as interpolation entered the reference: the discrepancy came out as exactly one step per episode boundary. The self-comparison worked as intended, but only one conclusion follows from it --- a duplicate diverges at every edit to the model. It has been deleted, the episode records are written by the reference itself, and there is nothing left to diverge. In passing the sample grew eightfold: 438 sets and 5323 episodes instead of 62 and 788. The second is the labelling of wakings: they were assigned to a cell by the \emph{maximum} damage during that same waking, so the label was partly created by the length being measured. The reverse causality is removed by labelling on the state at the \emph{beginning} of the waking.

% =====================================================================
\section{Three more traps, and the absence of an absolute null}\label{sec:null}
The list of statistical traps grew to ten, and the last three were found on our own rules rather than in the literature. Each of them had already produced a false verdict once before it was recognised.

\begin{enumerate}[start=8,leftmargin=1.9em,itemsep=3pt]
\item \textbf{A threshold the chosen statistic cannot reach.} A part of a rule required $p\le0.05$ from an exact permutation test over three scan points. That test's smallest attainable one-sided $p$ is $1/n!$, which at $n=3$ equals $0.167$: the threshold could not have been met by any behaviour of the model. The price is visible in the part that recorded $\rho=+1.000$ --- perfect monotonicity --- as ``not satisfied'' (section~\ref{sec:cap}). It is a relative of the sixth trap: both say that a criterion must be checked for \emph{feasibility} before the computation, not only for meaningfulness. The check is now built in: if the attainable $p$ exceeds the threshold, the part is declared unverifiable rather than failed.
\item \textbf{The nearest element of a grid that does not bracket.} A comparison against ``the same mean magnitude'' took the nearest grid point, and the grid ended before the value sought. The answer obtained was about the end of the grid. This trap is silent always, because a ``nearest'' element exists for any grid. The remedy is of two parts: the grid must bracket on both sides, and the agreement tolerance must be assigned before the computation from the construction of the scan --- half the smallest grid step --- so that failure to bracket makes the part unverifiable.
\item \textbf{The range of several noisy values is itself a random variable.} For $n$ independent values with relative error $r$, the expected range is a multiple of $r$ --- for three values, $1.693\,r$. A threshold ``the spread over $J$ is no more than 10\,\%'' is therefore meaningful only for $r<0.059$. At 300 sets $r=12\,\%$, so the rule was unverifiable; at 2000 sets it became verifiable --- and immediately failed, revealing the third route by which $J$ enters the model (section~\ref{sec:J}).
\end{enumerate}

\textbf{The central methodological result of the project is not a trap but a structural statement, and it applies to any model of this kind.} In a model with absorbing states and a heterogeneous population, \emph{there is no absolute null for any measure of accumulation}. A damaged cell does not recover on its own; over a long enough record the accumulation measure drifts upwards in the open model too, so ``no effect'' does not correspond to zero. Three consequences follow, and all three were paid for.

\begin{itemize}
\item A rule that tests an accumulation measure against zero cannot fail: the only question is how long the record is. The verdict it issues is about the length of the run.
\item Worse, the converse: a \emph{false} null passes not by being true but by lack of power. The rule reports ``satisfied'' and the defect is invisible from inside, because a null that cannot be violated is never violated.
\item Hence the construction that replaced it: every part of the rule is a \emph{matched difference}. The loop against a static equivalent levelled on the first half of the same record, paired within a set, on the same sid. The null of a matched difference is zero by construction, whatever the absorbing states do, because whatever they do they do on both sides.
\end{itemize}

\begin{ineq}
The general form of the statement. Let $D(t)$ be a measure of accumulation in a system with an absorbing set $A$ and a population heterogeneous in a parameter $v$. For every $v$ with positive hazard, $\mathrm{P}(x(t;v)\in A)\to1$, so $\mathrm{E}\,D(t)$ is non-decreasing and its limit is set by the fraction of the population with positive hazard, not by the mechanism under test. A test of $\mathrm{E}\,D(T)>0$ therefore has power tending to one against \emph{any} alternative, including the null one, as $T$ grows. The only quantities whose null is a fixed number are differences between two arms sharing the same absorbing dynamics --- which is why the comparison, and not the measurement, is the unit of evidence here.
\end{ineq}

% =====================================================================
\section{What the model predicts that is measurable}\label{sec:pred}
All of the work is done \emph{in silico}, and the first question put to it is what can be measured from here. Below is everything the model asserts about quantities accessible to experiment: a number, a direction, and what exactly refutes the assertion. The predictions are ordered by the strength of the test, not by the importance of the mechanism.

\begin{note}
A caveat about units, and it is obligatory. What is comparable is \emph{directions, orders of magnitude and ratios}, not absolute values. The parameters of the redox core are taken from the kinetics of cell lines and not of PV cells; the quantities of the sleep switch are expressed in millivolts of its own calibration. Where an absolute number stands below, it is quoted as the size of an effect inside the model, and it may not be carried over to experiment without recalculation.
\end{note}

\begin{enumerate}[leftmargin=1.5em,itemsep=4pt]
\item \textbf{Non-monotonicity of damage in vulnerability: the rescued band.} The population lives in three phases, not two. Between the healthy core and the collapsed one lies a band of cells which are \emph{more} vulnerable than their collapsed neighbours and are nonetheless intact: a fraction $0.00732\pm0.00042$ (seventeen standard errors from zero), and about $8.5\,\%$ of the population on a fine grid. With the slow defence frozen the band becomes exactly zero; under priming it widens $1.82$-fold. \emph{To measure}: damage to PV cells must not be a monotone function of their initial vulnerability --- the tissue must contain a subgroup of cells with high vulnerability markers and preserved parvalbumin, adjacent to damaged ones. Proxies for vulnerability: integrity of the perineuronal net, mitochondrial markers, indices of the oxidative load on the cell. \emph{Refuted by}: a monotone ``more vulnerable, more damaged'' relation at single-cell level. This is the most distinctive prediction of the work: it is counter-intuitive, follows from no averaged picture, and requires single-cell rather than tissue resolution.

\item \textbf{A precursor in sleep, and its source.} Sleep deteriorates \emph{before} the collapse threshold, and what does it is the slow memory of the defence rather than slow noise in the load: when the slow memory is taken away the height of the precursor falls in 107 sets of 174, $p=0.0015$. The height in the last day before a collapse is $1.061$ of the set's own healthy level, against an analytic bracket for the switch of $1.0500$--$1.1280$. \emph{To measure}: a lengthening of waking of the order of a few per cent in the days preceding an episode. Data of this kind are plentiful --- actigraphy and polysomnography in clinical high-risk cohorts and in schizophrenia, where sleep disturbance before an exacerbation is described. \emph{Refuted by}: the absence of a systematic shift before an episode; or by the presence of a shift that does not disappear under an intervention suppressing the slow link of the defence. The second is the stronger test but requires an animal model.

\item \textbf{The ceiling of passive repair, and the affine duration law.} The duration of an episode is set by the \emph{memory} of the defence and not by its strength: within a set, $99.8\,\%$ of the variance of the duration is explained by the course of the load and $0.2\,\%$ by repair, the decomposition being exact. The law is affine and does not saturate, duration $=1.28+0.97\,\tau_{\rm d}$ days, and a saturating form with a common offset is rejected. Hence: \emph{passive strengthening of repair does not attain the clinical target} --- missing by a factor of 2 to 10, since the required memory is about 12\,h whereas every measured value for the second stage lies in 24--88\,h. \emph{To measure}: interventions that lengthen or shorten the lifetime of the protective protein must shift the episode duration \emph{linearly}, without saturation; interventions that raise the strength of repair at unchanged memory must not shift it at all. \emph{Refuted by}: saturation of the response at large $\tau_{\rm d}$, or an appreciable shortening of episodes from strengthened repair at unchanged memory.

\item \textbf{Priming works through the top of the window, not through the speed of repair.} A preliminary mild load shortens the subsequent episode: the one-sided 95\,\% bound on the median duration is $1.51$ days, inside the clinical two. The mechanism is named: priming raises the \emph{top} of the window rather than accelerating repair. Quantitatively, accelerating the first stage of induction gives $-1$ to $-5\,\%$ of the duration, whereas strengthening it gives $-52\,\%$. \emph{To measure}: a hormesis protocol --- a mild oxidative exposure before the main one must shorten the episode. The two mechanisms are separated by an intervention that accelerates induction against one that raises its ceiling; the model predicts that only the second works. \emph{Refuted by}: equal efficacy of the two interventions, or an advantage for acceleration.

\item \textbf{The capacity loop: episodes rarer but longer.} Per-cell loss of defence capacity does not make collapses more frequent but makes them longer: duration $13.1\to25.6$ days with the rate falling $0.042\to0.033$ per set-day at risk, while episodicity is preserved at $0.75$ of the open model. Progression exceeds the capacity-matched static by $+0.0231$ at an MDE of $0.0096$, replicated on three independent noise realisations. \emph{To measure}: in longitudinal cohorts, as damage accumulates the episodes should become \emph{rarer and longer}, not more frequent. This is testable on data about the number and duration of exacerbations that are already being collected. \emph{Refuted by}: a growing rate of exacerbations under progression. Caveat: the absolute value of $25.6$ days lies far above the clinical two, so the working point of the loop is overstated; what is comparable here is the direction and the ratio.

\item \textbf{Hysteresis: retention depends on the depth to which the load is removed.} The broken state is retained after the load is removed, and the retention grows with the softness of the removal monotonically ($\rho=+1.000$, permutation $p=0.0083$): under the softest removal it is $0.561$ against $0.146$ for the open model. \emph{To measure}: the probability of relapse after treatment is withdrawn should depend monotonically on the completeness of the withdrawal rather than being constant. Testable in discontinuation studies, though ``completeness of removal'' needs an operationalisation of its own.

\item \textbf{A bursty load does not strike harder than a smooth one.} With the mean and the variance pinned, the collapse rate \emph{falls} by $40\,\%$ as the correlation time of the load grows. Rare strong pushes are no more dangerous than a steady rise --- a hypothesis we tested and refuted. \emph{To measure}: stress protocols with identical mean and variance but different temporal structure; clustered stress should give \emph{fewer} episodes than uniform stress. Testable in an animal model.

\item \textbf{A serotonin surge is a volume knob, not a trigger.} A surge is followed by $26\,\%$ of collapses, against an analytic expectation of $43\,\%$ and a pre-registered threshold of $80\,\%$; and no property of the episode changes. \emph{To measure}: serotonergic provocation should not reliably launch an episode and should not change its duration or depth.

\item \textbf{The asymmetry of the defence sets its mean level.} If the rise and decay constants of the protective protein differ by a factor $r$, the mean level of the defence stands at $w(r)=fr/(fr+1-f)$ with $f=0.705$, the measured fraction of waking. A scan of $r=0.33$--$3$ predicts $w=0.441$--$0.878$, and the measured levels of the defence follow that ordering exactly ($\rho=+1$, $p=0.0083$). The formula contains not one fitted number. \emph{To measure}: the two constants of the protective protein separately, in the cell. The model then predicts the mean level of the defence with no free parameters.

\item \textbf{Two negative predictions.} \emph{A loop through the common load cannot give chronicity and episodicity at once}: irreversibility engages at a loop strength of $1.11$ and the episodes die out at $0.18$ --- a sixfold structural gap, robust to a softer protocol. Consequence for experiment: if patients reliably show \emph{both} irreversible accumulation \emph{and} preserved episodicity, the mechanism does not lie in a loop through the common load. \emph{Breakdown of the rhythm as a locking staircase is unattainable at physiological coupling}: the first step requires $k\!\cdot\!D=4.77$\,mV while identically $k\!\cdot\!D\le k$, so at the working $k=2$\,mV it is not reached for any damage. The model thus predicts that sleep cycles of 48 and 72\,h should not occur in schizophrenia --- which agrees with their not being described. The prediction is weak, but it is the only place in the work where a disagreement between model and observation turned out to be in the model's favour.
\end{enumerate}

\textbf{What the model does \emph{not} predict.} Quantities of its internal construction --- the ratio of the window's boundaries, the positions of the folds in millivolts, the position of the load within the window, the number of cells in a set --- are not predictions: they are either chosen or auxiliary. Dopamine and the positive symptoms are outside the model by design. And, most importantly, the model describes the onset and the ending of an \emph{episode}, not the course of the disease, so any claim on the scale of years lies outside its domain of applicability.

% =====================================================================
\section{What we hold to be beyond dispute}\label{sec:solid}
The ledger of verdicts now contains 128 parts of rules: 66 satisfied, 48 not satisfied, 13 unverifiable, one below the bracket. Below is the part we are prepared to defend without reservation: claims with a threshold assigned before the computation from expectations and variances, with a paired control and with a check on refutability.

\textbf{On the mechanism.}
\begin{enumerate}[leftmargin=1.3em,itemsep=2pt,topsep=2pt]
\item The episode duration is a \emph{linear} function of the memory of the defence, without saturation; a test point outside the fit confirmed the form. Hence: passive strengthening of repair does not attain the clinical target, missing it by a factor of 2 to 10.
\item The spread of episode durations within a set is explained 99.8\,\% by the course of the load and 0.2\,\% by repair. The decomposition is exact by construction; the shares sum to within $10^{-16}$.
\item A long memory of the defence produces the precursor: when it is taken away the height falls in 107 sets of 174, $p=0.0015$. Slow noise in the load is not needed for this --- its justification has been withdrawn twice.
\item The asymmetry of the rise and decay constants is a regime parameter and not a detail: the mean level of the defence follows the derivation $w(r)$ with exact agreement of the ordering, $p=0.0083$ (section~\ref{sec:asym}).
\item The population lives in three phases. The rescued band exists (seventeen sigmas), becomes exactly \emph{zero} with the defence frozen, and widens twofold under priming (section~\ref{sec:band}).
\item The rhythm breaks as a locking staircase with rational rotation numbers and with alternating cycle lengths on the fractional steps; the cascade proceeds by adding one. The staircase is generated by branch switching of a map that does not cover the circle, not by mode locking, and in the whole window there is neither chaos nor a smooth invariant torus: the Lyapunov exponent along the orbit lies between $-3.54$ and $-1.62$ (section~\ref{sec:anal}).
\item A loop through the common load is a latch and not a knob: it releases 0.24\,\% of the broken sets against 1.000 for the open model, and it breaks sets where the static load does not suffice at all. It cannot give chronicity and episodicity at once --- the latch engages at a strength of 1.11 and the episodes die out at 0.18, a sixfold structural gap (section~\ref{sec:dopa}).
\item A loop acting per cell and on the capacity of the defence preserves episodicity at 0.75 of the open model and makes the episodes rarer and longer; progression exceeds the capacity-matched static by $+0.0231$ at an MDE of $0.0096$, replicated on three independent noise realisations (section~\ref{sec:cap}).
\end{enumerate}

\textbf{On method.}
\begin{enumerate}[leftmargin=1.3em,itemsep=2pt,topsep=2pt,resume]
\item The number of cells enters the model as an effective spread of resilience and not as a resolution of the computation; the limit $J\to\infty$ exists and has been measured.
\item Ten statistical traps, each of which would once have produced a false conclusion, and each closed by a check in the code.
\item In a model with absorbing states and a heterogeneous population there is no absolute null for any measure of accumulation; a false null there passes by lack of power rather than by being true, so the unit of evidence is a matched difference and not a measurement (section~\ref{sec:null}).
\item The sleep switch admits a closed-form reduction to a relay with hysteresis, confirmed independently three times --- on the folds (1.4503 and 2.4635 against 1.450 and 2.463), on the tongue edge (4.7651 against 4.77) and on the rotation numbers --- with the 10\,\% offset of the relay limit stated rather than concealed.
\item The identity of the accelerated computation is proved at two levels --- bitwise by fields and by all verdicts at once.
\end{enumerate}

\textbf{What we do \emph{not} claim.} The form of the dependence of the band width on the damaged fraction --- refuted twice with opposite signs. The explanation of the discrepancies in the hazard arithmetic by lag of the defence --- not confirmed in the form recorded; the reformulation in terms of the collapse rate looks sound but has no verdict yet and awaits a run on an independent configuration. The absolute duration produced by the capacity loop (25.6 days) lies far above the clinical two, so only its direction and ratio are claimed, not its magnitude; and the time constant of the dopamine loop is not taken from the literature, because no quantitative value for it exists there. Breakdown of the rhythm at the \emph{calibrated} ``damage--sleep'' coupling is unattainable, and the fact that the model reaches the staircase at a coupling twice as strong says nothing in itself about the disease. The quantitative form of $w(r)$, the thirteen parts with the status ``unverifiable'', and all quantities whose thresholds are finer than 3\,\% on counting measures, lie outside the circle of the indisputable.

% =====================================================================
\section{Open questions}\label{sec:crit}
\begin{enumerate}
\item \textbf{One cell instead of a population.} Damage to one cell transfers the load onto its neighbours; this may both spread the damage and contain it.
\item \textbf{Parameters from the wrong cells.} The kinetics of peroxiredoxin were measured in cell lines. The conclusion ``the cells stand at the edge of the window'' applies to them and not to PV cells.
\item \textbf{What the load is.} Part of the influx of oxidants is constant (the mitochondria), part depends on receptor activity and arrives from neighbouring cells. The shares are unknown.
\item \textbf{No memory between exposures.} \emph{Closed in version~4:} memory appeared twice --- as the two-stage defence with a constant of 63\,h and as priming (sections~\ref{sec:slow} and~\ref{sec:prime}). There is now something with which to describe the two-dose ketamine experiment.
\item \textbf{The daily rhythm: imposed or intrinsic.} Either the clock periodically changes the load, or the switch oscillates by itself and is merely entrained. In the second case the disease is an arrest of the cycle in the damaged phase.
\item \textbf{A phase shift greater than a quarter of a period} requires one more slow stage or an element of fractional order.
\item \textbf{The spectrum of slow fluctuations of the load} on scales of minutes and hours has not been measured, and whether noise can carry the system over depends on it.
\item \textbf{Death or silence.} The data on patients speak rather of silence of the cells.
\item \textbf{No development, no dopamine, no symptoms.} The model describes the state of a cell, not the disease.
\item \textbf{Which link is primary.} Oxidative defence, the BDNF support signal, or both as consequences of a common cause. What is to be counted as the slow variable and which times are to be substituted depend on the answer.
\item \textbf{The third slow quantity.} \emph{Half closed in version~5:} the asymmetry of the rise and decay of the defence is now in the model, has been measured and behaves according to the derivation $w(r)$ (section~\ref{sec:asym}), so this candidate is no longer a hypothesis. The second remains: the sleep precursor is explained by the memory of the defence only to a fifth, and there are no times in the equations other than those already built in. Either this is the filling of the population (which has no $\tau$ of its own), or one more slow stage is needed.
\item \textbf{What sets the intercept of the duration law.} The straight line $1.28+0.97\,\tau_{\rm d}$ gives 1.3 days at zero memory, and this remainder cannot be explained by the return of $x$: it returns within hours.
\item \textbf{Breakdown of the rhythm is unattainable at the working coupling.} The locking staircase exists, and the model \emph{does} reach it --- but at a ``damage--sleep'' coupling of no less than 4.77\,mV, that is, 2.4 times stronger than the working 2\,mV. The second way out recorded in version~4 (chronic damage holding the broken fraction high) is \emph{closed by measurement}: identically $kD\le k$, and in the broken sets the fraction is already at its ceiling. One of two things remains: either the coupling in the model is too weak, or the breakdown of the rhythm in the disease is not brought about by this route.
\item \textbf{Which loop closes the circle.} \emph{Half closed:} a loop through the common load is refuted as a model of the disease --- it gives either reversible episodes or irreversible shutdown, never both, and the cause is structural, since the load is one quantity for the whole population (section~\ref{sec:dopa}). A loop acting per cell, through the capacity of the defence, preserves episodicity and is the surviving candidate; but its working point produces episodes of 25.6 days, so its strength is not calibrated, and one part of its rule --- the course of the hazard over the matched static --- did not pass.
\item \textbf{Invariant tori and chaos in the switch.} \emph{Closed, negatively:} neither exists in the window, and the proof is an inequality on the exponent along the orbit rather than a scan (section~\ref{sec:anal}). What remains formally open is an exceptional parameter carrying a Denjoy-type Cantor minimal attractor; the argument against is structural, since such an attractor also requires the exponent at zero, and the exponent nowhere rises above $-1.6$. Also open: the overshoot exponent converged to $0.59$--$0.64$ against the theoretical $2/3$, so either a further decade in the fast time scale is needed or the correction contains a logarithm.
\item \textbf{The very idea of two states.} There is no direct experiment showing hysteresis in a PV cell; the two-dose ketamine experiment is the nearest observation to it.
\item \textbf{What sets the width of the rescued band.} The geometry of the boundary layer is refuted: two scans give opposite signs (section~\ref{sec:band}). So the width is determined not by the observable damaged fraction but by something third --- most probably the time the population has spent right at the tipping point. This is testable against the already measured fraction of time in the middle of the scale, and it must agree on \emph{both} scans, otherwise it is withdrawn in the same way.
\item \textbf{Why a short exposure strikes harder than a long one.} Lag of the defence was not confirmed in the form recorded (section~\ref{sec:burst}); the reformulation in terms of the collapse rate looks sound but has no verdict. If it too fails, the next candidate is a nonlinearity beyond the top of the load window: a surge goes past it whereas the daily range does not, and the hazard law there may be different.
\item \textbf{What the truth about counting quantities costs.} The collapse rate at the working step is underestimated by about 2\,\%, and this is a limit of the sampling, not of rounding (section~\ref{sec:prec}). Either the step is halved and twice as expensive, or the sub-step excursions are reconstructed exactly from the crossing fractions, which in the fused kernel already lie in registers.
\end{enumerate}

% =====================================================================
\section{Glossary}
{\small
\begin{description}[leftmargin=1.2em,itemsep=1pt,parsep=0pt,font=\normalfont\bfseries]
\item[AMPA receptor] the fast glutamate receptor
\item[ATP] the molecule in which a cell stores energy
\item[Astrocyte] a supporting cell of the brain; supplies neurons, including with the raw material for glutathione
\item[Axon] the output process of a neuron
\item[BDNF] brain-derived neurotrophic factor; a support signal released by active cells
\item[Catatonia] a state with severe motor disturbance; a separate diagnosis in ICD-11
\item[Cell culture] cells grown outside an organism
\item[Dendrite] an input process of a neuron
\item[Disinhibition] a rise in activity due to weakened inhibition
\item[Dopamine] a transmitter associated with the appraisal of the significance of events
\item[Enzyme] a protein that accelerates a chemical reaction
\item[Expression] production of a protein from a gene
\item[GABA] the main inhibitory transmitter
\item[GAD67] the enzyme that produces GABA
\item[Gamma rhythm] oscillations of cortical activity at 30--80\,Hz
\item[Glutamate] the main excitatory transmitter
\item[Glutathione] the principal consumable antioxidant of the cell
\item[Hippocampus] a brain region linked with memory; in Grace's scheme, the source of the excess signal
\item[ICD-11] the latest revision of the International Classification of Diseases
\item[Interneuron] an inhibitory neuron of the cortex
\item[Ketamine] a substance that blocks the NMDA receptor
\item[Master clock] the brain region above the crossing of the optic nerves (the suprachiasmatic nucleus), which synchronises the body's daily rhythms
\item[Membrane] the envelope of a cell
\item[Mitochondrion] the cell's power station; produces oxidants as a by-product
\item[mRNA] the working copy of a gene from which a protein is built
\item[NMDA receptor] the slow glutamate receptor; passes calcium and serves as an activity sensor
\item[Nox2] an enzyme that produces oxidants deliberately
\item[Oxidant] a molecule that takes electrons from other molecules
\item[Oxidative stress] an excess of oxidants over the capacity to clear them
\item[Parvalbumin (PV)] a calcium-binding protein; the marker of the fast inhibitory cells
\item[Perineuronal net] a shell of extracellular material around a PV cell
\item[Peroxiredoxin] the principal peroxide-scavenging enzyme
\item[Prodrome] the period before the first episode of psychosis
\item[Protein] a chain of amino acids assembled from a gene; the cell's main working material
\item[Psychosis] a state with hallucinations and delusions
\item[Pyramidal cell] an excitatory neuron of the cortex
\item[Receptor] a protein that receives a signal; in neurons, a gate for ions opened by a transmitter
\item[Serotonin] a transmitter governing the overall operating regime of the cortex; hallucinogens act on its receptors
\item[Sulfiredoxin] the enzyme that repairs hyperoxidised peroxiredoxin
\item[Synapse] the point of contact between two neurons
\item[Transmitter] a molecule that passes a signal across a synapse
\item[TrkB] the receptor that receives the BDNF signal
\end{description}}

% =====================================================================
\section{Sources}
{\small
\textbf{Verified while preparing the brochure} (tag \tP):
\begin{itemize}
\item Saha~S. et al. A systematic review of the prevalence of schizophrenia. PLoS Medicine, 2005.
\item Hilker~R. et al. Heritability of schizophrenia and schizophrenia spectrum based on the nationwide Danish twin register. Biological Psychiatry, 2018 (as reported in the press release).
\item Krystal~J.\,H. et al. Subanesthetic effects of the noncompetitive NMDA antagonist, ketamine, in humans. Archives of General Psychiatry, 1994.
\item Kapur~S. et al. Relationship between dopamine D2 occupancy, clinical response, and side effects. American Journal of Psychiatry, 2000.
\item Howes~O.\,D. et al. The nature of dopamine dysfunction in schizophrenia and what this means for treatment. Archives of General Psychiatry, 2012.
\item A study of glutathione turnover in the mouse brain, PubMed record 9165485.
\item Harkness~J.\,H. et al. Diurnal changes in perineuronal nets and parvalbumin neurons in the rat medial prefrontal cortex. bioRxiv preprint, 2020; published 2021.
\item Molecular and cellular rhythms in excitatory and inhibitory neurons in the mouse prefrontal cortex. bioRxiv preprint, 2024.
\item Schizophrenia and catatonia: from ICD-10 to ICD-11. Review article, Springer, 2025 (authors not verified). The list of ICD-11 codes of group 6A2 is from the code reference.
\item Hashimoto~T. et al. Relationship of brain-derived neurotrophic factor and its receptor TrkB to altered inhibitory prefrontal circuitry in schizophrenia. Journal of Neuroscience, 2005.
\item Fernandes~B.\,S. et al. Peripheral brain-derived neurotrophic factor in schizophrenia and the role of antipsychotics: meta-analysis and implications. Molecular Psychiatry, 2015.
\item Durany~N., Thome~J. Neurotrophic factors and the pathophysiology of schizophrenic psychoses. European Psychiatry, 2004; Angelucci~F. et al. BDNF in schizophrenia, depression and corresponding animal models. Molecular Psychiatry, 2005.
\item Jaggar~M., Vaidya~V.\,A. 5-HT2A receptors and BDNF regulation: implications for psychopathology. Book chapter, Springer, 2018 (from the abstract).
\item Papadia~S. et al. Synaptic NMDA receptor activity boosts intrinsic antioxidant defenses. Nature Neuroscience, 2008 (from a secondary account).
\end{itemize}
\textbf{Read by the authors of the project} (tag \tS):
\begin{itemize}
\item Hardingham~G.\,E., Do~K.\,Q. Linking early-life NMDAR hypofunction and oxidative stress in schizophrenia pathogenesis. Nature Reviews Neuroscience, 2016.
\item Cuenod~M. et al. Review of the vicious circles through oxidative stress, 2021, and the reply to criticism (Molecular Psychiatry, 2022).
\item Selvaggio~G., Coelho~P., Salvador~A. A model of the peroxiredoxin--thioredoxin system. Redox Biology, 2018.
\item Volman~V., Behrens~M.\,M., Sejnowski~T.\,J. A network model with loss of parvalbumin and GAD67, 2011.
\item Kinney~J.\,W. et al., 2006 --- ketamine in culture; Behrens~M.\,M. et al., 2008 --- ketamine in mice; Baxter~P.\,S. et al., 2015 --- the receptor and glutathione synthesis; Brennan~A.\,M. et al., 2009 --- the receptor and Nox2.
\item Kil~I.\,S. et al. Circadian oscillation of sulfiredoxin in the mitochondria. Molecular Cell, 2015; del~Olmo~M. et al. A robust model for circadian redox oscillations. International Journal of Molecular Sciences, 2019.
\item Edgar~R.\,S. et al., 2012 --- daily cycles of the peroxiredoxins; Wang~T.\,A. et al., 2012 --- oxidative state and excitability in the master clock; Kinoshita~C. et al., 2014 --- the daily rhythm of glutathione.
\item Harkness~J.\,H. et al. --- sleep disruption and oxidative stress in PV cells; Vyazovskiy~V.\,V. et al., 2009 --- firing rate and sleep; He~B.\,J., 2014 --- the scale-free spectrum of brain activity.
\item Work of A.~Grace's group on the MAM model; work of O.~Howes on dopamine synthesis before psychosis; Do~K.\,Q. et al., 2000, and the 2021 synthesis on glutathione in patients.
\item A review of the structure and function of PV cells; work on the stability of the cortex (Tsodyks, 1997; Ozeki, 2009); studies of the duration of symptoms and of the prodrome.
\end{itemize}
The bibliographic details of the second group are given from working notes and require verification before citation.

\textbf{Textbook knowledge} (tag \tU) is written from memory. Before citing such a fact it should be checked against a textbook, for example: Kandel~E. et al., Principles of Neural Science; Purves~D. et al., Neuroscience; Alberts~B. et al., Molecular Biology of the Cell.
}


\end{document}
